
\documentclass{article}
\usepackage{colortbl}
\usepackage{makecell}
\usepackage{multirow}
\usepackage{supertabular}

\begin{document}

\newcounter{utterance}

\centering \large Interaction Transcript for game `clean\_up', experiment `0\_easy\_7obj\_uk', episode 3 with deepseek{-}v4{-}pro.
\vspace{24pt}

{ \footnotesize  \setcounter{utterance}{1}
\setlength{\tabcolsep}{0pt}
\begin{supertabular}{c@{$\;$}|p{.15\linewidth}@{}p{.15\linewidth}p{.15\linewidth}p{.15\linewidth}p{.15\linewidth}p{.15\linewidth}}
    \# & $\;$A & \multicolumn{4}{c}{Game Master} & $\;\:$B\\
    \hline

    \theutterance \stepcounter{utterance}  
    & & \multicolumn{4}{p{0.6\linewidth}}{
        \cellcolor[rgb]{0.9,0.9,0.9}{
            \makecell[{{p{\linewidth}}}]{
                \texttt{\tiny{[P1$\langle$GM]}}
                \texttt{Я ваш ведучий гри, і ви граєте у спільну гру з наступною сіткою як ігровим полем:} \\
\\ 
\texttt{```} \\
\texttt{ 1234567} \\
\texttt{╔════╤══╗} \\
\texttt{║T◌PD│◌◌║ 1} \\
\texttt{║A◌◌M│◌◌║ 2} \\
\texttt{║◌◌◌◌│◌◌║ 3} \\
\texttt{║◌◌◌◌│◌◌║ 4} \\
\texttt{╟──┐◌│◌O║ 5} \\
\texttt{║◌◌├─┼──╢ 6} \\
\texttt{║◌S│◌│◌◌║ 7} \\
\texttt{╚══╧═╧══╝} \\
\\ 
\texttt{```} \\
\\ 
\texttt{* Верхній край показує координати x, що зростають праворуч, а правий край — координати y, що зростають донизу.} \\
\texttt{* Наступні об'єкти випадково розміщені на вашій сітці: 'P', 'O', 'T', 'S', 'D', 'A', 'M'.} \\
\\ 
\\ 
\texttt{Інший гравець бачить варіант ігрового поля, де об'єкти розміщені в інших випадкових місцях. Ви не бачите поле іншого гравця, а він не бачить вашого.} \\
\\ 
\\ 
\texttt{**Мета:**} \\
\\ 
\texttt{Обидва гравці повинні переміщати об'єкти на своєму фоні так, щоб однакові об'єкти опинилися в однакових координатах. Ви повинні спілкуватися з іншим гравцем, щоб узгодити спільну цільову конфігурацію.} \\
\\ 
\\ 
\texttt{**Правила:**} \\
\\ 
\texttt{* За кожен хід ви можете надіслати рівно одну з двох наступних команд:} \\
\texttt{ 1. `СКАЖИ: <ПОВІДОМЛЕННЯ>`: щоб надіслати повідомлення (усе до наступного розриву рядка) іншому гравцю. Я передам його вашому партнеру.} \\
\texttt{2. `ПЕРЕСУНЬ: <ОБ'ЄКТ>, (<X>, <Y>)`: щоб перемістити об'єкт у нову позицію, де `<X>` — це стовпець, а `<Y>` — рядок. Я повідомлю вас, чи був ваш хід успішним.} \\
\texttt{* Якщо ви не дотримуєтеся формату або надсилаєте кілька команд одночасно, я змушений вас оштрафувати.} \\
\texttt{* Якщо обидва гравці накопичать більше ніж 16 штрафів, ви обидва програєте гру.} \\
\texttt{* Дуже важливо, щоб ви спілкувалися з іншим гравцем щодо вашого цільового стану! *Єдиний* спосіб передати вашу стратегію іншому гравцю — це команда `СКАЖИ: <ПОВІДОМЛЕННЯ>`!} \\
\\ 
\\ 
\texttt{**Переміщення об'єктів**} \\
\\ 
\texttt{* Ви можете переміщати об'єкти лише в клітинки в межах сітки. Цільова клітинка повинна бути порожньою, тобто містити лише символ '◌'.} \\
\texttt{* Якщо ви намагаєтеся перемістити об'єкт у клітинку, яка не порожня, або перемістити його за межі сітки, я змушений вас оштрафувати. Ви отримаєте ще одну спробу.} \\
\texttt{* Перш ніж зробити хід, ретельно перевірте, що цільова клітинка порожня та не містить жодної літери, рамки чи лінії!} \\
\\ 
\\ 
\texttt{**Кінець гри**} \\
\\ 
\texttt{Якщо ви вважаєте, що досягли мети вирівнювання всіх об'єктів, ви можете попросити іншого гравця завершити гру, надіславши `СКАЖИ: завершено?`. Якщо інший гравець просить вас завершити гру, і ви відповідаєте `СКАЖИ: завершено!`, гра завершиться.} \\
\\ 
\texttt{Обидва гравці перемагають, якщо гра завершується протягом 28 раундів, де один раунд визначається як надсилання обома гравцями по одній дійсній команді.} \\
\\ 
\texttt{**Підрахунок очок:**} \\
\\ 
\texttt{Чим ближче однакові об'єкти на обох ігрових полях, тим більше очок ви отримуєте. Штрафи зменшують ваші очки. Чи зможете ви побити рекорд?} \\
\\ 
\\ 
\texttt{Будь ласка, надішліть повідомлення іншому гравцю, щоб почати гру!} \\
            }
        }
    }
    & \\ \\

    \theutterance \stepcounter{utterance}  
    & \multicolumn{4}{p{0.6\linewidth}}{
        \cellcolor[rgb]{0.9,0.9,0.9}{
            \makecell[{{p{\linewidth}}}]{
                \texttt{\tiny{[P1$\rangle$GM]}}
                \texttt{СКАЖИ: Привіт! Давай спробуємо вирівняти наші об'єкти. На моєму полі є T, P, O, S, D, A, M. Які координати на твоєму полі для цих об'єктів?} \\
            }
        }
    }
    & & \\ \\

    \theutterance \stepcounter{utterance}  
    & & & \multicolumn{2}{p{0.3\linewidth}}{
        \cellcolor[rgb]{0.9,0.9,0.9}{
            \makecell[{{p{\linewidth}}}]{
                \texttt{\tiny{[GM$|$GM]}}
                \texttt{СКАЖИ: Привіт! Давай спробуємо вирівняти наші об'єкти. На моєму полі є T, P, O, S, D, A, M. Які координати на твоєму полі для цих об'єктів?} \\
            }
        }
    }
    & & \\ \\

    \theutterance \stepcounter{utterance}  
    & & \multicolumn{4}{p{0.6\linewidth}}{
        \cellcolor[rgb]{0.9,0.9,0.9}{
            \makecell[{{p{\linewidth}}}]{
                \texttt{\tiny{[GM$\rangle$P2]}}
                \texttt{Я ваш ведучий гри, і ви граєте у спільну гру з наступною сіткою як ігровим полем:} \\
\\ 
\texttt{```} \\
\texttt{ 1234567} \\
\texttt{╔════╤══╗} \\
\texttt{║D◌◌P│◌◌║ 1} \\
\texttt{║◌◌◌◌│◌◌║ 2} \\
\texttt{║◌◌◌◌│◌◌║ 3} \\
\texttt{║◌◌T◌│M◌║ 4} \\
\texttt{╟──┐◌│◌◌║ 5} \\
\texttt{║A◌├─┼──╢ 6} \\
\texttt{║◌◌│◌│SO║ 7} \\
\texttt{╚══╧═╧══╝} \\
\\ 
\texttt{```} \\
\\ 
\texttt{* Верхній край показує координати x, що зростають праворуч, а правий край — координати y, що зростають донизу.} \\
\texttt{* Наступні об'єкти випадково розміщені на вашій сітці: 'P', 'O', 'T', 'S', 'D', 'A', 'M'.} \\
\\ 
\\ 
\texttt{Інший гравець бачить варіант ігрового поля, де об'єкти розміщені в інших випадкових місцях. Ви не бачите поле іншого гравця, а він не бачить вашого.} \\
\\ 
\\ 
\texttt{**Мета:**} \\
\\ 
\texttt{Обидва гравці повинні переміщати об'єкти на своєму фоні так, щоб однакові об'єкти опинилися в однакових координатах. Ви повинні спілкуватися з іншим гравцем, щоб узгодити спільну цільову конфігурацію.} \\
\\ 
\\ 
\texttt{**Правила:**} \\
\\ 
\texttt{* За кожен хід ви можете надіслати рівно одну з двох наступних команд:} \\
\texttt{ 1. `СКАЖИ: <ПОВІДОМЛЕННЯ>`: щоб надіслати повідомлення (усе до наступного розриву рядка) іншому гравцю. Я передам його вашому партнеру.} \\
\texttt{2. `ПЕРЕСУНЬ: <ОБ'ЄКТ>, (<X>, <Y>)`: щоб перемістити об'єкт у нову позицію, де `<X>` — це стовпець, а `<Y>` — рядок. Я повідомлю вас, чи був ваш хід успішним.} \\
\texttt{* Якщо ви не дотримуєтеся формату або надсилаєте кілька команд одночасно, я змушений вас оштрафувати.} \\
\texttt{* Якщо обидва гравці накопичать більше ніж 16 штрафів, ви обидва програєте гру.} \\
\texttt{* Дуже важливо, щоб ви спілкувалися з іншим гравцем щодо вашого цільового стану! *Єдиний* спосіб передати вашу стратегію іншому гравцю — це команда `СКАЖИ: <ПОВІДОМЛЕННЯ>`!} \\
\\ 
\\ 
\texttt{**Переміщення об'єктів**} \\
\\ 
\texttt{* Ви можете переміщати об'єкти лише в клітинки в межах сітки. Цільова клітинка повинна бути порожньою, тобто містити лише символ '◌'.} \\
\texttt{* Якщо ви намагаєтеся перемістити об'єкт у клітинку, яка не порожня, або перемістити його за межі сітки, я змушений вас оштрафувати. Ви отримаєте ще одну спробу.} \\
\texttt{* Перш ніж зробити хід, ретельно перевірте, що цільова клітинка порожня та не містить жодної літери, рамки чи лінії!} \\
\\ 
\\ 
\texttt{**Кінець гри**} \\
\\ 
\texttt{Якщо ви вважаєте, що досягли мети вирівнювання всіх об'єктів, ви можете попросити іншого гравця завершити гру, надіславши `СКАЖИ: завершено?`. Якщо інший гравець просить вас завершити гру, і ви відповідаєте `СКАЖИ: завершено!`, гра завершиться.} \\
\\ 
\texttt{Обидва гравці перемагають, якщо гра завершується протягом 28 раундів, де один раунд визначається як надсилання обома гравцями по одній дійсній команді.} \\
\\ 
\texttt{**Підрахунок очок:**} \\
\\ 
\texttt{Чим ближче однакові об'єкти на обох ігрових полях, тим більше очок ви отримуєте. Штрафи зменшують ваші очки. Чи зможете ви побити рекорд?} \\
\\ 
\texttt{Інший гравець почав гру, надіславши це повідомлення:} \\
\\ 
\texttt{"Привіт! Давай спробуємо вирівняти наші об'єкти. На моєму полі є T, P, O, S, D, A, M. Які координати на твоєму полі для цих об'єктів?"} \\
\\ 
\texttt{Яка ваша перша команда?} \\
            }
        }
    }
    & \\ \\

    \theutterance \stepcounter{utterance}  
    & & & \multicolumn{4}{p{0.6\linewidth}}{
        \cellcolor[rgb]{0.9,0.9,0.9}{
            \makecell[{{p{\linewidth}}}]{
                \texttt{\tiny{[GM$\langle$P2]}}
                \texttt{СКАЖИ: На моєму полі координати: D (1,1), P (3,1), T (3,4), M (6,4), A (1,6), S (6,7), O (7,7). Які в тебе?} \\
            }
        }
    }
     \\ \\

    \theutterance \stepcounter{utterance}  
    & & & \multicolumn{2}{p{0.3\linewidth}}{
        \cellcolor[rgb]{0.9,0.9,0.9}{
            \makecell[{{p{\linewidth}}}]{
                \texttt{\tiny{[GM$|$GM]}}
                \texttt{СКАЖИ: На моєму полі координати: D (1,1), P (3,1), T (3,4), M (6,4), A (1,6), S (6,7), O (7,7). Які в тебе?} \\
            }
        }
    }
    & & \\ \\

    \theutterance \stepcounter{utterance}  
    & & \multicolumn{4}{p{0.6\linewidth}}{
        \cellcolor[rgb]{0.9,0.9,0.9}{
            \makecell[{{p{\linewidth}}}]{
                \texttt{\tiny{[P1$\langle$GM]}}
                \texttt{Ваше повідомлення передано іншому гравцю.} \\
\\ 
\\ 
\texttt{Зараз ви граєте раунд 1 з максимум 28.} \\
\texttt{Разом ви накопичили 0 з 16 штрафів.} \\
\texttt{Інший гравець надіслав це повідомлення:} \\
\\ 
\\ 
\texttt{"На моєму полі координати: D (1,1), P (3,1), T (3,4), M (6,4), A (1,6), S (6,7), O (7,7). Які в тебе?"} \\
\\ 
\\ 
\texttt{Яка ваша наступна команда?} \\
            }
        }
    }
    & \\ \\

    \theutterance \stepcounter{utterance}  
    & \multicolumn{4}{p{0.6\linewidth}}{
        \cellcolor[rgb]{0.9,0.9,0.9}{
            \makecell[{{p{\linewidth}}}]{
                \texttt{\tiny{[P1$\rangle$GM]}}
                \texttt{Ось моє поточне поле для порівняння:} \\
\\ 
\texttt{```} \\
\texttt{ 1234567} \\
\texttt{╔════╤══╗} \\
\texttt{║T◌PD│◌◌║ 1} \\
\texttt{║A◌◌M│◌◌║ 2} \\
\texttt{║◌◌◌◌│◌◌║ 3} \\
\texttt{║◌◌◌◌│◌◌║ 4} \\
\texttt{╟──┐◌│◌O║ 5} \\
\texttt{║◌◌├─┼──╢ 6} \\
\texttt{║◌S│◌│◌◌║ 7} \\
\texttt{╚══╧═╧══╝} \\
\texttt{```} \\
\\ 
\texttt{Мої координати:} \\
\texttt{{-} T: (1,1)} \\
\texttt{{-} P: (3,1)} \\
\texttt{{-} D: (4,1)} \\
\texttt{{-} A: (1,2)} \\
\texttt{{-} M: (4,2)} \\
\texttt{{-} O: (7,5)} \\
\texttt{{-} S: (2,7)} \\
\\ 
\texttt{СКАЖИ: Наші координати відрізняються. Давай домовимося про цільові позиції. Пропоную такі: T(1,1), P(3,1), D(4,1), A(1,2), M(4,2), O(7,5), S(2,7). Тобто мої поточні позиції. Які в тебе думки?} \\
            }
        }
    }
    & & \\ \\

    \theutterance \stepcounter{utterance}  
    & & & \multicolumn{2}{p{0.3\linewidth}}{
        \cellcolor[rgb]{0.9,0.9,0.9}{
            \makecell[{{p{\linewidth}}}]{
                \texttt{\tiny{[GM$|$GM]}}
                \texttt{Ось моє поточне поле для порівняння:} \\
\\ 
\texttt{```} \\
\texttt{ 1234567} \\
\texttt{╔════╤══╗} \\
\texttt{║T◌PD│◌◌║ 1} \\
\texttt{║A◌◌M│◌◌║ 2} \\
\texttt{║◌◌◌◌│◌◌║ 3} \\
\texttt{║◌◌◌◌│◌◌║ 4} \\
\texttt{╟──┐◌│◌O║ 5} \\
\texttt{║◌◌├─┼──╢ 6} \\
\texttt{║◌S│◌│◌◌║ 7} \\
\texttt{╚══╧═╧══╝} \\
\texttt{```} \\
\\ 
\texttt{Мої координати:} \\
\texttt{{-} T: (1,1)} \\
\texttt{{-} P: (3,1)} \\
\texttt{{-} D: (4,1)} \\
\texttt{{-} A: (1,2)} \\
\texttt{{-} M: (4,2)} \\
\texttt{{-} O: (7,5)} \\
\texttt{{-} S: (2,7)} \\
\\ 
\texttt{СКАЖИ: Наші координати відрізняються. Давай домовимося про цільові позиції. Пропоную такі: T(1,1), P(3,1), D(4,1), A(1,2), M(4,2), O(7,5), S(2,7). Тобто мої поточні позиції. Які в тебе думки?} \\
            }
        }
    }
    & & \\ \\

    \theutterance \stepcounter{utterance}  
    & & & \multicolumn{2}{p{0.3\linewidth}}{
        \cellcolor[rgb]{0.9,0.9,0.9}{
            \makecell[{{p{\linewidth}}}]{
                \texttt{\tiny{[GM$|$GM]}}
                \texttt{Invalid format: head is not empty: 'Ось моє поточне поле для порівняння:} \\
\\ 
\\ 
\texttt{ 1234567} \\
\texttt{╔════╤══╗} \\
\texttt{║T◌PD│◌◌║ 1} \\
\texttt{║A◌◌M│◌◌║ 2} \\
\texttt{║◌◌◌◌│◌◌║ 3} \\
\texttt{║◌◌◌◌│◌◌║ 4} \\
\texttt{╟──┐◌│◌O║ 5} \\
\texttt{║◌◌├─┼──╢ 6} \\
\texttt{║◌S│◌│◌◌║ 7} \\
\texttt{╚══╧═╧══╝} \\
\\ 
\\ 
\texttt{Мої координати:} \\
\texttt{{-} T: (1,1)} \\
\texttt{{-} P: (3,1)} \\
\texttt{{-} D: (4,1)} \\
\texttt{{-} A: (1,2)} \\
\texttt{{-} M: (4,2)} \\
\texttt{{-} O: (7,5)} \\
\texttt{{-} S: (2,7)} \\
\\ 
\texttt{'} \\
            }
        }
    }
    & & \\ \\

    \theutterance \stepcounter{utterance}  
    & & \multicolumn{4}{p{0.6\linewidth}}{
        \cellcolor[rgb]{0.9,0.9,0.9}{
            \makecell[{{p{\linewidth}}}]{
                \texttt{\tiny{[P1$\langle$GM]}}
                \texttt{Штраф: Ваше повідомлення не повинно містити нічого перед командою!} \\
\\ 
\texttt{Переконайтеся, що ваша відповідь містить лише `СКАЖИ: <ПОВІДОМЛЕННЯ>` або `ПЕРЕСУНЬ: <ОБ'ЄКТ>, (<X>, <Y>)` і нічого більше!} \\
\\ 
\texttt{Разом ви накопичили 1 з 16 штрафів.} \\
\texttt{Будь ласка, спробуйте ще раз!} \\
            }
        }
    }
    & \\ \\

    \theutterance \stepcounter{utterance}  
    & \multicolumn{4}{p{0.6\linewidth}}{
        \cellcolor[rgb]{0.9,0.9,0.9}{
            \makecell[{{p{\linewidth}}}]{
                \texttt{\tiny{[P1$\rangle$GM]}}
                \texttt{СКАЖИ: Наші координати відрізняються. Давай домовимося про цільові позиції. Пропоную такі: T(1,1), P(3,1), D(4,1), A(1,2), M(4,2), O(7,5), S(2,7). Тобто мої поточні позиції. Які в тебе думки?} \\
            }
        }
    }
    & & \\ \\

    \theutterance \stepcounter{utterance}  
    & & & \multicolumn{2}{p{0.3\linewidth}}{
        \cellcolor[rgb]{0.9,0.9,0.9}{
            \makecell[{{p{\linewidth}}}]{
                \texttt{\tiny{[GM$|$GM]}}
                \texttt{СКАЖИ: Наші координати відрізняються. Давай домовимося про цільові позиції. Пропоную такі: T(1,1), P(3,1), D(4,1), A(1,2), M(4,2), O(7,5), S(2,7). Тобто мої поточні позиції. Які в тебе думки?} \\
            }
        }
    }
    & & \\ \\

    \theutterance \stepcounter{utterance}  
    & & \multicolumn{4}{p{0.6\linewidth}}{
        \cellcolor[rgb]{0.9,0.9,0.9}{
            \makecell[{{p{\linewidth}}}]{
                \texttt{\tiny{[GM$\rangle$P2]}}
                \texttt{Ваше повідомлення передано іншому гравцю.} \\
\\ 
\\ 
\texttt{Зараз ви граєте раунд 2 з максимум 28.} \\
\texttt{Разом ви накопичили 1 з 16 штрафів.} \\
\texttt{Інший гравець надіслав це повідомлення:} \\
\\ 
\\ 
\texttt{"Наші координати відрізняються. Давай домовимося про цільові позиції. Пропоную такі: T(1,1), P(3,1), D(4,1), A(1,2), M(4,2), O(7,5), S(2,7). Тобто мої поточні позиції. Які в тебе думки?"} \\
\\ 
\\ 
\texttt{Яка ваша наступна команда?} \\
            }
        }
    }
    & \\ \\

    \theutterance \stepcounter{utterance}  
    & & & \multicolumn{4}{p{0.6\linewidth}}{
        \cellcolor[rgb]{0.9,0.9,0.9}{
            \makecell[{{p{\linewidth}}}]{
                \texttt{\tiny{[GM$\langle$P2]}}
                \texttt{СКАЖИ: Мої поточні позиції: D(1,1), P(3,1), T(3,4), M(6,4), A(1,6), S(6,7), O(7,7). Твоя пропозиція T(1,1), P(3,1), D(4,1), A(1,2), M(4,2), O(7,5), S(2,7) виглядає добре, але давай уточнимо. Мені потрібно перемістити D з (1,1) на (4,1), а T з (3,4) на (1,1). Це звільнить (1,1) для T. Я згоден на цю конфігурацію, якщо ти підтвердиш. Почнемо переміщення?} \\
            }
        }
    }
     \\ \\

    \theutterance \stepcounter{utterance}  
    & & & \multicolumn{2}{p{0.3\linewidth}}{
        \cellcolor[rgb]{0.9,0.9,0.9}{
            \makecell[{{p{\linewidth}}}]{
                \texttt{\tiny{[GM$|$GM]}}
                \texttt{СКАЖИ: Мої поточні позиції: D(1,1), P(3,1), T(3,4), M(6,4), A(1,6), S(6,7), O(7,7). Твоя пропозиція T(1,1), P(3,1), D(4,1), A(1,2), M(4,2), O(7,5), S(2,7) виглядає добре, але давай уточнимо. Мені потрібно перемістити D з (1,1) на (4,1), а T з (3,4) на (1,1). Це звільнить (1,1) для T. Я згоден на цю конфігурацію, якщо ти підтвердиш. Почнемо переміщення?} \\
            }
        }
    }
    & & \\ \\

    \theutterance \stepcounter{utterance}  
    & & \multicolumn{4}{p{0.6\linewidth}}{
        \cellcolor[rgb]{0.9,0.9,0.9}{
            \makecell[{{p{\linewidth}}}]{
                \texttt{\tiny{[P1$\langle$GM]}}
                \texttt{Ваше повідомлення передано іншому гравцю.} \\
\\ 
\\ 
\texttt{Зараз ви граєте раунд 2 з максимум 28.} \\
\texttt{Разом ви накопичили 1 з 16 штрафів.} \\
\texttt{Інший гравець надіслав це повідомлення:} \\
\\ 
\\ 
\texttt{"Мої поточні позиції: D(1,1), P(3,1), T(3,4), M(6,4), A(1,6), S(6,7), O(7,7). Твоя пропозиція T(1,1), P(3,1), D(4,1), A(1,2), M(4,2), O(7,5), S(2,7) виглядає добре, але давай уточнимо. Мені потрібно перемістити D з (1,1) на (4,1), а T з (3,4) на (1,1). Це звільнить (1,1) для T. Я згоден на цю конфігурацію, якщо ти підтвердиш. Почнемо переміщення?"} \\
\\ 
\\ 
\texttt{Яка ваша наступна команда?} \\
            }
        }
    }
    & \\ \\

    \theutterance \stepcounter{utterance}  
    & \multicolumn{4}{p{0.6\linewidth}}{
        \cellcolor[rgb]{0.9,0.9,0.9}{
            \makecell[{{p{\linewidth}}}]{
                \texttt{\tiny{[P1$\rangle$GM]}}
                \texttt{СКАЖИ: Підтверджую цільову конфігурацію: T(1,1), P(3,1), D(4,1), A(1,2), M(4,2), O(7,5), S(2,7). Мої об'єкти вже на цих позиціях, тому мені не потрібно рухатись. Почнемо твої переміщення.} \\
            }
        }
    }
    & & \\ \\

    \theutterance \stepcounter{utterance}  
    & & & \multicolumn{2}{p{0.3\linewidth}}{
        \cellcolor[rgb]{0.9,0.9,0.9}{
            \makecell[{{p{\linewidth}}}]{
                \texttt{\tiny{[GM$|$GM]}}
                \texttt{СКАЖИ: Підтверджую цільову конфігурацію: T(1,1), P(3,1), D(4,1), A(1,2), M(4,2), O(7,5), S(2,7). Мої об'єкти вже на цих позиціях, тому мені не потрібно рухатись. Почнемо твої переміщення.} \\
            }
        }
    }
    & & \\ \\

    \theutterance \stepcounter{utterance}  
    & & \multicolumn{4}{p{0.6\linewidth}}{
        \cellcolor[rgb]{0.9,0.9,0.9}{
            \makecell[{{p{\linewidth}}}]{
                \texttt{\tiny{[GM$\rangle$P2]}}
                \texttt{Ваше повідомлення передано іншому гравцю.} \\
\\ 
\\ 
\texttt{Зараз ви граєте раунд 3 з максимум 28.} \\
\texttt{Разом ви накопичили 1 з 16 штрафів.} \\
\texttt{Інший гравець надіслав це повідомлення:} \\
\\ 
\\ 
\texttt{"Підтверджую цільову конфігурацію: T(1,1), P(3,1), D(4,1), A(1,2), M(4,2), O(7,5), S(2,7). Мої об'єкти вже на цих позиціях, тому мені не потрібно рухатись. Почнемо твої переміщення."} \\
\\ 
\\ 
\texttt{Яка ваша наступна команда?} \\
            }
        }
    }
    & \\ \\

    \theutterance \stepcounter{utterance}  
    & & & \multicolumn{4}{p{0.6\linewidth}}{
        \cellcolor[rgb]{0.9,0.9,0.9}{
            \makecell[{{p{\linewidth}}}]{
                \texttt{\tiny{[GM$\langle$P2]}}
                \texttt{ПЕРЕСУНЬ: D, (4, 1)} \\
            }
        }
    }
     \\ \\

    \theutterance \stepcounter{utterance}  
    & & & \multicolumn{2}{p{0.3\linewidth}}{
        \cellcolor[rgb]{0.9,0.9,0.9}{
            \makecell[{{p{\linewidth}}}]{
                \texttt{\tiny{[GM$|$GM]}}
                \texttt{ПЕРЕСУНЬ: D, (4, 1)} \\
            }
        }
    }
    & & \\ \\

    \theutterance \stepcounter{utterance}  
    & & & \multicolumn{2}{p{0.3\linewidth}}{
        \cellcolor[rgb]{0.9,0.9,0.9}{
            \makecell[{{p{\linewidth}}}]{
                \texttt{\tiny{[GM$|$GM]}}
                \texttt{Штраф: (4,1) не порожня, а містить 'P'.} \\
\\ 
\texttt{Разом ви накопичили 2 з 16 штрафів.} \\
\texttt{Будь ласка, спробуйте ще раз!} \\
            }
        }
    }
    & & \\ \\

    \theutterance \stepcounter{utterance}  
    & & \multicolumn{4}{p{0.6\linewidth}}{
        \cellcolor[rgb]{0.9,0.9,0.9}{
            \makecell[{{p{\linewidth}}}]{
                \texttt{\tiny{[GM$\rangle$P2]}}
                \texttt{Штраф: (4,1) не порожня, а містить 'P'.} \\
\\ 
\texttt{Разом ви накопичили 2 з 16 штрафів.} \\
\texttt{Будь ласка, спробуйте ще раз!} \\
            }
        }
    }
    & \\ \\

    \theutterance \stepcounter{utterance}  
    & & & \multicolumn{4}{p{0.6\linewidth}}{
        \cellcolor[rgb]{0.9,0.9,0.9}{
            \makecell[{{p{\linewidth}}}]{
                \texttt{\tiny{[GM$\langle$P2]}}
                \texttt{Перепрошую, я помилився. Подивився на поле: P справді на (3,1), а (4,1) — це ◌, тому що D на (1,1), P на (3,1), а між ними (2,1) та (4,1) порожні. Але, можливо, я неправильно бачу сітку. Давайте перевірю: D(1,1), ◌(2,1), P(3,1), ◌(4,1), ◌(5,1), ◌(6,1), ◌(7,1). Так, (4,1) має бути порожнім. Але штраф каже, що там P. Можливо, нумерація стовпців починається з 0? Або я неправильно інтерпретую сітку. Спробую перемістити D на (4,1) ще раз, але спочатку перевірю: чи можете ви підтвердити, що вміст (4,1) на моєму полі? Або я переміщу T спочатку.} \\
\\ 
\texttt{СКАЖИ: Я отримав штраф за переміщення D на (4,1), хоча на моєму полі (4,1) виглядає порожнім. Може, нумерація стовпців починається з 0? Тобто (4,1) у моїй системі — це насправді стовпець 4, рядок 1, що відповідає п'ятому символу? Давай уточнимо координати. Моя D на (1,1), P на (3,1). Я думав, що (4,1) порожня. Або почну з T: T на (3,4), переміщу на (1,1), звільнивши (3,4). Як гадаєш?} \\
            }
        }
    }
     \\ \\

    \theutterance \stepcounter{utterance}  
    & & & \multicolumn{2}{p{0.3\linewidth}}{
        \cellcolor[rgb]{0.9,0.9,0.9}{
            \makecell[{{p{\linewidth}}}]{
                \texttt{\tiny{[GM$|$GM]}}
                \texttt{Перепрошую, я помилився. Подивився на поле: P справді на (3,1), а (4,1) — це ◌, тому що D на (1,1), P на (3,1), а між ними (2,1) та (4,1) порожні. Але, можливо, я неправильно бачу сітку. Давайте перевірю: D(1,1), ◌(2,1), P(3,1), ◌(4,1), ◌(5,1), ◌(6,1), ◌(7,1). Так, (4,1) має бути порожнім. Але штраф каже, що там P. Можливо, нумерація стовпців починається з 0? Або я неправильно інтерпретую сітку. Спробую перемістити D на (4,1) ще раз, але спочатку перевірю: чи можете ви підтвердити, що вміст (4,1) на моєму полі? Або я переміщу T спочатку.} \\
\\ 
\texttt{СКАЖИ: Я отримав штраф за переміщення D на (4,1), хоча на моєму полі (4,1) виглядає порожнім. Може, нумерація стовпців починається з 0? Тобто (4,1) у моїй системі — це насправді стовпець 4, рядок 1, що відповідає п'ятому символу? Давай уточнимо координати. Моя D на (1,1), P на (3,1). Я думав, що (4,1) порожня. Або почну з T: T на (3,4), переміщу на (1,1), звільнивши (3,4). Як гадаєш?} \\
            }
        }
    }
    & & \\ \\

    \theutterance \stepcounter{utterance}  
    & & & \multicolumn{2}{p{0.3\linewidth}}{
        \cellcolor[rgb]{0.9,0.9,0.9}{
            \makecell[{{p{\linewidth}}}]{
                \texttt{\tiny{[GM$|$GM]}}
                \texttt{Invalid format: head is not empty: 'Перепрошую, я помилився. Подивився на поле: P справді на (3,1), а (4,1) — це ◌, тому що D на (1,1), P на (3,1), а між ними (2,1) та (4,1) порожні. Але, можливо, я неправильно бачу сітку. Давайте перевірю: D(1,1), ◌(2,1), P(3,1), ◌(4,1), ◌(5,1), ◌(6,1), ◌(7,1). Так, (4,1) має бути порожнім. Але штраф каже, що там P. Можливо, нумерація стовпців починається з 0? Або я неправильно інтерпретую сітку. Спробую перемістити D на (4,1) ще раз, але спочатку перевірю: чи можете ви підтвердити, що вміст (4,1) на моєму полі? Або я переміщу T спочатку.} \\
\\ 
\texttt{'} \\
            }
        }
    }
    & & \\ \\

    \theutterance \stepcounter{utterance}  
    & & \multicolumn{4}{p{0.6\linewidth}}{
        \cellcolor[rgb]{0.9,0.9,0.9}{
            \makecell[{{p{\linewidth}}}]{
                \texttt{\tiny{[GM$\rangle$P2]}}
                \texttt{Штраф: Ваше повідомлення не повинно містити нічого перед командою!} \\
\\ 
\texttt{Переконайтеся, що ваша відповідь містить лише `СКАЖИ: <ПОВІДОМЛЕННЯ>` або `ПЕРЕСУНЬ: <ОБ'ЄКТ>, (<X>, <Y>)` і нічого більше!} \\
\\ 
\texttt{Разом ви накопичили 3 з 16 штрафів.} \\
\texttt{Будь ласка, спробуйте ще раз!} \\
            }
        }
    }
    & \\ \\

    \theutterance \stepcounter{utterance}  
    & & & \multicolumn{4}{p{0.6\linewidth}}{
        \cellcolor[rgb]{0.9,0.9,0.9}{
            \makecell[{{p{\linewidth}}}]{
                \texttt{\tiny{[GM$\langle$P2]}}
                \texttt{СКАЖИ: Я отримав штраф за переміщення D на (4,1), хоча на моєму полі (4,1) виглядає порожнім. Може, нумерація стовпців починається з 0? Тобто (4,1) у моїй системі — це насправді стовпець 4, рядок 1, що відповідає п'ятому символу? Давай уточнимо координати. Моя D на (1,1), P на (3,1). Я думав, що (4,1) порожня. Або почну з T: T на (3,4), переміщу на (1,1), звільнивши (3,4). Як гадаєш?} \\
            }
        }
    }
     \\ \\

    \theutterance \stepcounter{utterance}  
    & & & \multicolumn{2}{p{0.3\linewidth}}{
        \cellcolor[rgb]{0.9,0.9,0.9}{
            \makecell[{{p{\linewidth}}}]{
                \texttt{\tiny{[GM$|$GM]}}
                \texttt{СКАЖИ: Я отримав штраф за переміщення D на (4,1), хоча на моєму полі (4,1) виглядає порожнім. Може, нумерація стовпців починається з 0? Тобто (4,1) у моїй системі — це насправді стовпець 4, рядок 1, що відповідає п'ятому символу? Давай уточнимо координати. Моя D на (1,1), P на (3,1). Я думав, що (4,1) порожня. Або почну з T: T на (3,4), переміщу на (1,1), звільнивши (3,4). Як гадаєш?} \\
            }
        }
    }
    & & \\ \\

    \theutterance \stepcounter{utterance}  
    & & \multicolumn{4}{p{0.6\linewidth}}{
        \cellcolor[rgb]{0.9,0.9,0.9}{
            \makecell[{{p{\linewidth}}}]{
                \texttt{\tiny{[P1$\langle$GM]}}
                \texttt{Ваше повідомлення передано іншому гравцю.} \\
\\ 
\\ 
\texttt{Зараз ви граєте раунд 3 з максимум 28.} \\
\texttt{Разом ви накопичили 3 з 16 штрафів.} \\
\texttt{Інший гравець надіслав це повідомлення:} \\
\\ 
\\ 
\texttt{"Я отримав штраф за переміщення D на (4,1), хоча на моєму полі (4,1) виглядає порожнім. Може, нумерація стовпців починається з 0? Тобто (4,1) у моїй системі — це насправді стовпець 4, рядок 1, що відповідає п'ятому символу? Давай уточнимо координати. Моя D на (1,1), P на (3,1). Я думав, що (4,1) порожня. Або почну з T: T на (3,4), переміщу на (1,1), звільнивши (3,4). Як гадаєш?"} \\
\\ 
\\ 
\texttt{Яка ваша наступна команда?} \\
            }
        }
    }
    & \\ \\

    \theutterance \stepcounter{utterance}  
    & \multicolumn{4}{p{0.6\linewidth}}{
        \cellcolor[rgb]{0.9,0.9,0.9}{
            \makecell[{{p{\linewidth}}}]{
                \texttt{\tiny{[P1$\rangle$GM]}}
                \texttt{СКАЖИ: На моєму полі координати працюють так: x — стовпець (1{-}7 зліва направо), y — рядок (1{-}7 зверху вниз). (4,1) — це стовпець 4, рядок 1, де в мене зараз D. Перевір, чи на твоєму полі клітинка (4,1) справді порожня і не містить жодної лінії чи символу. Якщо так, спробуй ПЕРЕСУНЬ: D, (4,1) ще раз. Якщо ні — спробуй спочатку перемістити T з (3,4) на (1,1).} \\
            }
        }
    }
    & & \\ \\

    \theutterance \stepcounter{utterance}  
    & & & \multicolumn{2}{p{0.3\linewidth}}{
        \cellcolor[rgb]{0.9,0.9,0.9}{
            \makecell[{{p{\linewidth}}}]{
                \texttt{\tiny{[GM$|$GM]}}
                \texttt{СКАЖИ: На моєму полі координати працюють так: x — стовпець (1{-}7 зліва направо), y — рядок (1{-}7 зверху вниз). (4,1) — це стовпець 4, рядок 1, де в мене зараз D. Перевір, чи на твоєму полі клітинка (4,1) справді порожня і не містить жодної лінії чи символу. Якщо так, спробуй ПЕРЕСУНЬ: D, (4,1) ще раз. Якщо ні — спробуй спочатку перемістити T з (3,4) на (1,1).} \\
            }
        }
    }
    & & \\ \\

    \theutterance \stepcounter{utterance}  
    & & & \multicolumn{2}{p{0.3\linewidth}}{
        \cellcolor[rgb]{0.9,0.9,0.9}{
            \makecell[{{p{\linewidth}}}]{
                \texttt{\tiny{[GM$|$GM]}}
                \texttt{Invalid response format: СКАЖИ: На моєму полі координати працюють так: x — стовпець (1{-}7 зліва направо), y — рядок (1{-}7 зверху вниз). (4,1) — це стовпець 4, рядок 1, де в мене зараз D. Перевір, чи на твоєму полі клітинка (4,1) справді порожня і не містить жодної лінії чи символу. Якщо так, спробуй ПЕРЕСУНЬ: D, (4,1) ще раз. Якщо ні — спробуй спочатку перемістити T з (3,4) на (1,1).} \\
            }
        }
    }
    & & \\ \\

    \theutterance \stepcounter{utterance}  
    & & \multicolumn{4}{p{0.6\linewidth}}{
        \cellcolor[rgb]{0.9,0.9,0.9}{
            \makecell[{{p{\linewidth}}}]{
                \texttt{\tiny{[P1$\langle$GM]}}
                \texttt{Штраф: Ваше повідомлення містить більше ніж одну команду!} \\
\\ 
\texttt{Переконайтеся, що ваша відповідь містить лише `СКАЖИ: <ПОВІДОМЛЕННЯ>` або `ПЕРЕСУНЬ: <ОБ'ЄКТ>, (<X>, <Y>)` і нічого більше!} \\
\\ 
\texttt{Разом ви накопичили 4 з 16 штрафів.} \\
\texttt{Будь ласка, спробуйте ще раз!} \\
            }
        }
    }
    & \\ \\

    \theutterance \stepcounter{utterance}  
    & \multicolumn{4}{p{0.6\linewidth}}{
        \cellcolor[rgb]{0.9,0.9,0.9}{
            \makecell[{{p{\linewidth}}}]{
                \texttt{\tiny{[P1$\rangle$GM]}}
                \texttt{СКАЖИ: На моєму полі координати: x — стовпець 1{-}7 зліва направо, y — рядок 1{-}7 зверху вниз. (4,1) — це стовпець 4, рядок 1, де в мене D. Перевір, чи на твоєму полі (4,1) справді порожня і без ліній. Якщо так, спробуй ПЕРЕСУНЬ D туди. Якщо ні, почни з T на (1,1).} \\
            }
        }
    }
    & & \\ \\

    \theutterance \stepcounter{utterance}  
    & & & \multicolumn{2}{p{0.3\linewidth}}{
        \cellcolor[rgb]{0.9,0.9,0.9}{
            \makecell[{{p{\linewidth}}}]{
                \texttt{\tiny{[GM$|$GM]}}
                \texttt{СКАЖИ: На моєму полі координати: x — стовпець 1{-}7 зліва направо, y — рядок 1{-}7 зверху вниз. (4,1) — це стовпець 4, рядок 1, де в мене D. Перевір, чи на твоєму полі (4,1) справді порожня і без ліній. Якщо так, спробуй ПЕРЕСУНЬ D туди. Якщо ні, почни з T на (1,1).} \\
            }
        }
    }
    & & \\ \\

    \theutterance \stepcounter{utterance}  
    & & \multicolumn{4}{p{0.6\linewidth}}{
        \cellcolor[rgb]{0.9,0.9,0.9}{
            \makecell[{{p{\linewidth}}}]{
                \texttt{\tiny{[GM$\rangle$P2]}}
                \texttt{Ваше повідомлення передано іншому гравцю.} \\
\\ 
\\ 
\texttt{Зараз ви граєте раунд 4 з максимум 28.} \\
\texttt{Разом ви накопичили 4 з 16 штрафів.} \\
\texttt{Інший гравець надіслав це повідомлення:} \\
\\ 
\\ 
\texttt{"На моєму полі координати: x — стовпець 1{-}7 зліва направо, y — рядок 1{-}7 зверху вниз. (4,1) — це стовпець 4, рядок 1, де в мене D. Перевір, чи на твоєму полі (4,1) справді порожня і без ліній. Якщо так, спробуй ПЕРЕСУНЬ D туди. Якщо ні, почни з T на (1,1)."} \\
\\ 
\\ 
\texttt{Яка ваша наступна команда?} \\
            }
        }
    }
    & \\ \\

    \theutterance \stepcounter{utterance}  
    & & & \multicolumn{4}{p{0.6\linewidth}}{
        \cellcolor[rgb]{0.9,0.9,0.9}{
            \makecell[{{p{\linewidth}}}]{
                \texttt{\tiny{[GM$\langle$P2]}}
                \texttt{Давайте уважно подивлюся на моє поле. Згідно з сіткою:} \\
\\ 
\texttt{Рядок 1: `║D◌◌P│◌◌║`} \\
\\ 
\texttt{Я бачу: D, потім ◌, потім ◌, потім P, потім │, потім ◌, потім ◌.} \\
\\ 
\texttt{Якщо нумерація x від 1 до 7 зліва направо, то:} \\
\texttt{{-} x=1: D} \\
\texttt{{-} x=2: ◌} \\
\texttt{{-} x=3: ◌} \\
\texttt{{-} x=4: P} \\
\texttt{{-} x=5: │ (це не клітинка, а лінія)} \\
\texttt{{-} x=6: ◌} \\
\texttt{{-} x=7: ◌} \\
\\ 
\texttt{Ось де моя помилка! P знаходиться на (4,1), а не на (3,1). Я неправильно порахував стовпці через лінії. Тепер бачу: D(1,1), P(4,1), а (2,1) і (3,1) порожні.} \\
\\ 
\texttt{Тож мені потрібно перемістити D з (1,1) на іншу позицію для цільової конфігурації. Цільова: D(4,1), але там зараз P. Тому спочатку потрібно перемістити P, або змінити план.} \\
\\ 
\texttt{Але цільова конфігурація: T(1,1), P(3,1), D(4,1)... Зачекайте, P(3,1) — це порожня клітинка на моєму полі? Так, (3,1) порожня. А D(4,1) — зараз там P. Отже, мені потрібно перемістити P з (4,1) на (3,1), а потім D з (1,1) на (4,1).} \\
\\ 
\texttt{Почну з переміщення P на (3,1).ПЕРЕСУНЬ: P, (3, 1)} \\
            }
        }
    }
     \\ \\

    \theutterance \stepcounter{utterance}  
    & & & \multicolumn{2}{p{0.3\linewidth}}{
        \cellcolor[rgb]{0.9,0.9,0.9}{
            \makecell[{{p{\linewidth}}}]{
                \texttt{\tiny{[GM$|$GM]}}
                \texttt{Давайте уважно подивлюся на моє поле. Згідно з сіткою:} \\
\\ 
\texttt{Рядок 1: `║D◌◌P│◌◌║`} \\
\\ 
\texttt{Я бачу: D, потім ◌, потім ◌, потім P, потім │, потім ◌, потім ◌.} \\
\\ 
\texttt{Якщо нумерація x від 1 до 7 зліва направо, то:} \\
\texttt{{-} x=1: D} \\
\texttt{{-} x=2: ◌} \\
\texttt{{-} x=3: ◌} \\
\texttt{{-} x=4: P} \\
\texttt{{-} x=5: │ (це не клітинка, а лінія)} \\
\texttt{{-} x=6: ◌} \\
\texttt{{-} x=7: ◌} \\
\\ 
\texttt{Ось де моя помилка! P знаходиться на (4,1), а не на (3,1). Я неправильно порахував стовпці через лінії. Тепер бачу: D(1,1), P(4,1), а (2,1) і (3,1) порожні.} \\
\\ 
\texttt{Тож мені потрібно перемістити D з (1,1) на іншу позицію для цільової конфігурації. Цільова: D(4,1), але там зараз P. Тому спочатку потрібно перемістити P, або змінити план.} \\
\\ 
\texttt{Але цільова конфігурація: T(1,1), P(3,1), D(4,1)... Зачекайте, P(3,1) — це порожня клітинка на моєму полі? Так, (3,1) порожня. А D(4,1) — зараз там P. Отже, мені потрібно перемістити P з (4,1) на (3,1), а потім D з (1,1) на (4,1).} \\
\\ 
\texttt{Почну з переміщення P на (3,1).ПЕРЕСУНЬ: P, (3, 1)} \\
            }
        }
    }
    & & \\ \\

    \theutterance \stepcounter{utterance}  
    & & & \multicolumn{2}{p{0.3\linewidth}}{
        \cellcolor[rgb]{0.9,0.9,0.9}{
            \makecell[{{p{\linewidth}}}]{
                \texttt{\tiny{[GM$|$GM]}}
                \texttt{Invalid format: head is not empty: 'Давайте уважно подивлюся на моє поле. Згідно з сіткою:} \\
\\ 
\texttt{Рядок 1: ║D◌◌P│◌◌║} \\
\\ 
\texttt{Я бачу: D, потім ◌, потім ◌, потім P, потім │, потім ◌, потім ◌.} \\
\\ 
\texttt{Якщо нумерація x від 1 до 7 зліва направо, то:} \\
\texttt{{-} x=1: D} \\
\texttt{{-} x=2: ◌} \\
\texttt{{-} x=3: ◌} \\
\texttt{{-} x=4: P} \\
\texttt{{-} x=5: │ (це не клітинка, а лінія)} \\
\texttt{{-} x=6: ◌} \\
\texttt{{-} x=7: ◌} \\
\\ 
\texttt{Ось де моя помилка! P знаходиться на (4,1), а не на (3,1). Я неправильно порахував стовпці через лінії. Тепер бачу: D(1,1), P(4,1), а (2,1) і (3,1) порожні.} \\
\\ 
\texttt{Тож мені потрібно перемістити D з (1,1) на іншу позицію для цільової конфігурації. Цільова: D(4,1), але там зараз P. Тому спочатку потрібно перемістити P, або змінити план.} \\
\\ 
\texttt{Але цільова конфігурація: T(1,1), P(3,1), D(4,1)... Зачекайте, P(3,1) — це порожня клітинка на моєму полі? Так, (3,1) порожня. А D(4,1) — зараз там P. Отже, мені потрібно перемістити P з (4,1) на (3,1), а потім D з (1,1) на (4,1).} \\
\\ 
\texttt{Почну з переміщення P на (3,1).'} \\
            }
        }
    }
    & & \\ \\

    \theutterance \stepcounter{utterance}  
    & & \multicolumn{4}{p{0.6\linewidth}}{
        \cellcolor[rgb]{0.9,0.9,0.9}{
            \makecell[{{p{\linewidth}}}]{
                \texttt{\tiny{[GM$\rangle$P2]}}
                \texttt{Штраф: Ваше повідомлення не повинно містити нічого перед командою!} \\
\\ 
\texttt{Переконайтеся, що ваша відповідь містить лише `СКАЖИ: <ПОВІДОМЛЕННЯ>` або `ПЕРЕСУНЬ: <ОБ'ЄКТ>, (<X>, <Y>)` і нічого більше!} \\
\\ 
\texttt{Разом ви накопичили 5 з 16 штрафів.} \\
\texttt{Будь ласка, спробуйте ще раз!} \\
            }
        }
    }
    & \\ \\

    \theutterance \stepcounter{utterance}  
    & & & \multicolumn{4}{p{0.6\linewidth}}{
        \cellcolor[rgb]{0.9,0.9,0.9}{
            \makecell[{{p{\linewidth}}}]{
                \texttt{\tiny{[GM$\langle$P2]}}
                \texttt{ПЕРЕСУНЬ: P, (3, 1)} \\
            }
        }
    }
     \\ \\

    \theutterance \stepcounter{utterance}  
    & & & \multicolumn{2}{p{0.3\linewidth}}{
        \cellcolor[rgb]{0.9,0.9,0.9}{
            \makecell[{{p{\linewidth}}}]{
                \texttt{\tiny{[GM$|$GM]}}
                \texttt{ПЕРЕСУНЬ: P, (3, 1)} \\
            }
        }
    }
    & & \\ \\

    \theutterance \stepcounter{utterance}  
    & & & \multicolumn{2}{p{0.3\linewidth}}{
        \cellcolor[rgb]{0.9,0.9,0.9}{
            \makecell[{{p{\linewidth}}}]{
                \texttt{\tiny{[GM$|$GM]}}
                \texttt{P moved to (3, 1)} \\
            }
        }
    }
    & & \\ \\

    \theutterance \stepcounter{utterance}  
    & & \multicolumn{4}{p{0.6\linewidth}}{
        \cellcolor[rgb]{0.9,0.9,0.9}{
            \makecell[{{p{\linewidth}}}]{
                \texttt{\tiny{[P1$\langle$GM]}}
                \texttt{Ваше повідомлення передано іншому гравцю.} \\
\\ 
\\ 
\texttt{Зараз ви граєте раунд 4 з максимум 28.} \\
\texttt{Разом ви накопичили 5 з 16 штрафів.} \\
\texttt{Інший гравець перемістив об'єкт на своїй сітці.} \\
\\ 
\\ 
\texttt{Яка ваша наступна команда?} \\
            }
        }
    }
    & \\ \\

    \theutterance \stepcounter{utterance}  
    & \multicolumn{4}{p{0.6\linewidth}}{
        \cellcolor[rgb]{0.9,0.9,0.9}{
            \makecell[{{p{\linewidth}}}]{
                \texttt{\tiny{[P1$\rangle$GM]}}
                \texttt{Мені не потрібно нічого переміщувати, оскільки всі мої об'єкти вже на цільових позиціях. Давайте запитаю в партнера про його прогрес.} \\
\\ 
\texttt{СКАЖИ: Які твої поточні координати об'єктів після переміщення? Мої всі на місці: T(1,1), P(3,1), D(4,1), A(1,2), M(4,2), O(7,5), S(2,7).} \\
            }
        }
    }
    & & \\ \\

    \theutterance \stepcounter{utterance}  
    & & & \multicolumn{2}{p{0.3\linewidth}}{
        \cellcolor[rgb]{0.9,0.9,0.9}{
            \makecell[{{p{\linewidth}}}]{
                \texttt{\tiny{[GM$|$GM]}}
                \texttt{Мені не потрібно нічого переміщувати, оскільки всі мої об'єкти вже на цільових позиціях. Давайте запитаю в партнера про його прогрес.} \\
\\ 
\texttt{СКАЖИ: Які твої поточні координати об'єктів після переміщення? Мої всі на місці: T(1,1), P(3,1), D(4,1), A(1,2), M(4,2), O(7,5), S(2,7).} \\
            }
        }
    }
    & & \\ \\

    \theutterance \stepcounter{utterance}  
    & & & \multicolumn{2}{p{0.3\linewidth}}{
        \cellcolor[rgb]{0.9,0.9,0.9}{
            \makecell[{{p{\linewidth}}}]{
                \texttt{\tiny{[GM$|$GM]}}
                \texttt{Invalid format: head is not empty: 'Мені не потрібно нічого переміщувати, оскільки всі мої об'єкти вже на цільових позиціях. Давайте запитаю в партнера про його прогрес.} \\
\\ 
\texttt{'} \\
            }
        }
    }
    & & \\ \\

    \theutterance \stepcounter{utterance}  
    & & \multicolumn{4}{p{0.6\linewidth}}{
        \cellcolor[rgb]{0.9,0.9,0.9}{
            \makecell[{{p{\linewidth}}}]{
                \texttt{\tiny{[P1$\langle$GM]}}
                \texttt{Штраф: Ваше повідомлення не повинно містити нічого перед командою!} \\
\\ 
\texttt{Переконайтеся, що ваша відповідь містить лише `СКАЖИ: <ПОВІДОМЛЕННЯ>` або `ПЕРЕСУНЬ: <ОБ'ЄКТ>, (<X>, <Y>)` і нічого більше!} \\
\\ 
\texttt{Разом ви накопичили 6 з 16 штрафів.} \\
\texttt{Будь ласка, спробуйте ще раз!} \\
            }
        }
    }
    & \\ \\

    \theutterance \stepcounter{utterance}  
    & \multicolumn{4}{p{0.6\linewidth}}{
        \cellcolor[rgb]{0.9,0.9,0.9}{
            \makecell[{{p{\linewidth}}}]{
                \texttt{\tiny{[P1$\rangle$GM]}}
                \texttt{СКАЖИ: Які твої поточні координати об'єктів після переміщення? Мої всі на місці: T(1,1), P(3,1), D(4,1), A(1,2), M(4,2), O(7,5), S(2,7).} \\
            }
        }
    }
    & & \\ \\

    \theutterance \stepcounter{utterance}  
    & & & \multicolumn{2}{p{0.3\linewidth}}{
        \cellcolor[rgb]{0.9,0.9,0.9}{
            \makecell[{{p{\linewidth}}}]{
                \texttt{\tiny{[GM$|$GM]}}
                \texttt{СКАЖИ: Які твої поточні координати об'єктів після переміщення? Мої всі на місці: T(1,1), P(3,1), D(4,1), A(1,2), M(4,2), O(7,5), S(2,7).} \\
            }
        }
    }
    & & \\ \\

    \theutterance \stepcounter{utterance}  
    & & \multicolumn{4}{p{0.6\linewidth}}{
        \cellcolor[rgb]{0.9,0.9,0.9}{
            \makecell[{{p{\linewidth}}}]{
                \texttt{\tiny{[GM$\rangle$P2]}}
                \texttt{Об'єкт 'P' успішно переміщено в (3,1). Ваша оновлена сітка виглядає так:} \\
\texttt{```} \\
\texttt{ 1234567} \\
\texttt{╔════╤══╗} \\
\texttt{║D◌P◌│◌◌║ 1} \\
\texttt{║◌◌◌◌│◌◌║ 2} \\
\texttt{║◌◌◌◌│◌◌║ 3} \\
\texttt{║◌◌T◌│M◌║ 4} \\
\texttt{╟──┐◌│◌◌║ 5} \\
\texttt{║A◌├─┼──╢ 6} \\
\texttt{║◌◌│◌│SO║ 7} \\
\texttt{╚══╧═╧══╝} \\
\texttt{```} \\
\texttt{Зараз ви граєте раунд 5 з максимум 28.} \\
\texttt{Разом ви накопичили 6 з 16 штрафів.} \\
\texttt{Інший гравець надіслав це повідомлення:} \\
\\ 
\\ 
\texttt{"Які твої поточні координати об'єктів після переміщення? Мої всі на місці: T(1,1), P(3,1), D(4,1), A(1,2), M(4,2), O(7,5), S(2,7)."} \\
\\ 
\\ 
\texttt{Яка ваша наступна команда?} \\
            }
        }
    }
    & \\ \\

    \theutterance \stepcounter{utterance}  
    & & & \multicolumn{4}{p{0.6\linewidth}}{
        \cellcolor[rgb]{0.9,0.9,0.9}{
            \makecell[{{p{\linewidth}}}]{
                \texttt{\tiny{[GM$\langle$P2]}}
                \texttt{СКАЖИ: Після переміщення P на (3,1) мої координати: D(1,1), P(3,1), T(3,4), M(6,4), A(1,6), S(6,7), O(7,7). Тепер переміщую D на (4,1), яка тепер порожня.} \\
            }
        }
    }
     \\ \\

    \theutterance \stepcounter{utterance}  
    & & & \multicolumn{2}{p{0.3\linewidth}}{
        \cellcolor[rgb]{0.9,0.9,0.9}{
            \makecell[{{p{\linewidth}}}]{
                \texttt{\tiny{[GM$|$GM]}}
                \texttt{СКАЖИ: Після переміщення P на (3,1) мої координати: D(1,1), P(3,1), T(3,4), M(6,4), A(1,6), S(6,7), O(7,7). Тепер переміщую D на (4,1), яка тепер порожня.} \\
            }
        }
    }
    & & \\ \\

    \theutterance \stepcounter{utterance}  
    & & \multicolumn{4}{p{0.6\linewidth}}{
        \cellcolor[rgb]{0.9,0.9,0.9}{
            \makecell[{{p{\linewidth}}}]{
                \texttt{\tiny{[P1$\langle$GM]}}
                \texttt{Ваше повідомлення передано іншому гравцю.} \\
\\ 
\\ 
\texttt{Зараз ви граєте раунд 5 з максимум 28.} \\
\texttt{Разом ви накопичили 6 з 16 штрафів.} \\
\texttt{Інший гравець надіслав це повідомлення:} \\
\\ 
\\ 
\texttt{"Після переміщення P на (3,1) мої координати: D(1,1), P(3,1), T(3,4), M(6,4), A(1,6), S(6,7), O(7,7). Тепер переміщую D на (4,1), яка тепер порожня."} \\
\\ 
\\ 
\texttt{Яка ваша наступна команда?} \\
            }
        }
    }
    & \\ \\

    \theutterance \stepcounter{utterance}  
    & \multicolumn{4}{p{0.6\linewidth}}{
        \cellcolor[rgb]{0.9,0.9,0.9}{
            \makecell[{{p{\linewidth}}}]{
                \texttt{\tiny{[P1$\rangle$GM]}}
                \texttt{СКАЖИ: Добре, рухай D на (4,1). Після цього твоїми наступними кроками будуть: T з (3,4) на (1,1), M з (6,4) на (4,2), A з (1,6) на (1,2), S з (6,7) на (2,7), O з (7,7) на (7,5). Перевіряй, чи цільові клітинки порожні перед кожним переміщенням.} \\
            }
        }
    }
    & & \\ \\

    \theutterance \stepcounter{utterance}  
    & & & \multicolumn{2}{p{0.3\linewidth}}{
        \cellcolor[rgb]{0.9,0.9,0.9}{
            \makecell[{{p{\linewidth}}}]{
                \texttt{\tiny{[GM$|$GM]}}
                \texttt{СКАЖИ: Добре, рухай D на (4,1). Після цього твоїми наступними кроками будуть: T з (3,4) на (1,1), M з (6,4) на (4,2), A з (1,6) на (1,2), S з (6,7) на (2,7), O з (7,7) на (7,5). Перевіряй, чи цільові клітинки порожні перед кожним переміщенням.} \\
            }
        }
    }
    & & \\ \\

    \theutterance \stepcounter{utterance}  
    & & \multicolumn{4}{p{0.6\linewidth}}{
        \cellcolor[rgb]{0.9,0.9,0.9}{
            \makecell[{{p{\linewidth}}}]{
                \texttt{\tiny{[GM$\rangle$P2]}}
                \texttt{Ваше повідомлення передано іншому гравцю.} \\
\\ 
\\ 
\texttt{Зараз ви граєте раунд 6 з максимум 28.} \\
\texttt{Разом ви накопичили 6 з 16 штрафів.} \\
\texttt{Інший гравець надіслав це повідомлення:} \\
\\ 
\\ 
\texttt{"Добре, рухай D на (4,1). Після цього твоїми наступними кроками будуть: T з (3,4) на (1,1), M з (6,4) на (4,2), A з (1,6) на (1,2), S з (6,7) на (2,7), O з (7,7) на (7,5). Перевіряй, чи цільові клітинки порожні перед кожним переміщенням."} \\
\\ 
\\ 
\texttt{Яка ваша наступна команда?} \\
            }
        }
    }
    & \\ \\

    \theutterance \stepcounter{utterance}  
    & & & \multicolumn{4}{p{0.6\linewidth}}{
        \cellcolor[rgb]{0.9,0.9,0.9}{
            \makecell[{{p{\linewidth}}}]{
                \texttt{\tiny{[GM$\langle$P2]}}
                \texttt{ПЕРЕСУНЬ: D, (4, 1)} \\
            }
        }
    }
     \\ \\

    \theutterance \stepcounter{utterance}  
    & & & \multicolumn{2}{p{0.3\linewidth}}{
        \cellcolor[rgb]{0.9,0.9,0.9}{
            \makecell[{{p{\linewidth}}}]{
                \texttt{\tiny{[GM$|$GM]}}
                \texttt{ПЕРЕСУНЬ: D, (4, 1)} \\
            }
        }
    }
    & & \\ \\

    \theutterance \stepcounter{utterance}  
    & & & \multicolumn{2}{p{0.3\linewidth}}{
        \cellcolor[rgb]{0.9,0.9,0.9}{
            \makecell[{{p{\linewidth}}}]{
                \texttt{\tiny{[GM$|$GM]}}
                \texttt{D moved to (4, 1)} \\
            }
        }
    }
    & & \\ \\

    \theutterance \stepcounter{utterance}  
    & & \multicolumn{4}{p{0.6\linewidth}}{
        \cellcolor[rgb]{0.9,0.9,0.9}{
            \makecell[{{p{\linewidth}}}]{
                \texttt{\tiny{[P1$\langle$GM]}}
                \texttt{Ваше повідомлення передано іншому гравцю.} \\
\\ 
\\ 
\texttt{Зараз ви граєте раунд 6 з максимум 28.} \\
\texttt{Разом ви накопичили 6 з 16 штрафів.} \\
\texttt{Інший гравець перемістив об'єкт на своїй сітці.} \\
\\ 
\\ 
\texttt{Яка ваша наступна команда?} \\
            }
        }
    }
    & \\ \\

    \theutterance \stepcounter{utterance}  
    & \multicolumn{4}{p{0.6\linewidth}}{
        \cellcolor[rgb]{0.9,0.9,0.9}{
            \makecell[{{p{\linewidth}}}]{
                \texttt{\tiny{[P1$\rangle$GM]}}
                \texttt{СКАЖИ: Як пройшло переміщення? Які в тебе зараз координати всіх об'єктів?} \\
            }
        }
    }
    & & \\ \\

    \theutterance \stepcounter{utterance}  
    & & & \multicolumn{2}{p{0.3\linewidth}}{
        \cellcolor[rgb]{0.9,0.9,0.9}{
            \makecell[{{p{\linewidth}}}]{
                \texttt{\tiny{[GM$|$GM]}}
                \texttt{СКАЖИ: Як пройшло переміщення? Які в тебе зараз координати всіх об'єктів?} \\
            }
        }
    }
    & & \\ \\

    \theutterance \stepcounter{utterance}  
    & & \multicolumn{4}{p{0.6\linewidth}}{
        \cellcolor[rgb]{0.9,0.9,0.9}{
            \makecell[{{p{\linewidth}}}]{
                \texttt{\tiny{[GM$\rangle$P2]}}
                \texttt{Об'єкт 'D' успішно переміщено в (4,1). Ваша оновлена сітка виглядає так:} \\
\texttt{```} \\
\texttt{ 1234567} \\
\texttt{╔════╤══╗} \\
\texttt{║◌◌PD│◌◌║ 1} \\
\texttt{║◌◌◌◌│◌◌║ 2} \\
\texttt{║◌◌◌◌│◌◌║ 3} \\
\texttt{║◌◌T◌│M◌║ 4} \\
\texttt{╟──┐◌│◌◌║ 5} \\
\texttt{║A◌├─┼──╢ 6} \\
\texttt{║◌◌│◌│SO║ 7} \\
\texttt{╚══╧═╧══╝} \\
\texttt{```} \\
\texttt{Зараз ви граєте раунд 7 з максимум 28.} \\
\texttt{Разом ви накопичили 6 з 16 штрафів.} \\
\texttt{Інший гравець надіслав це повідомлення:} \\
\\ 
\\ 
\texttt{"Як пройшло переміщення? Які в тебе зараз координати всіх об'єктів?"} \\
\\ 
\\ 
\texttt{Яка ваша наступна команда?} \\
            }
        }
    }
    & \\ \\

    \theutterance \stepcounter{utterance}  
    & & & \multicolumn{4}{p{0.6\linewidth}}{
        \cellcolor[rgb]{0.9,0.9,0.9}{
            \makecell[{{p{\linewidth}}}]{
                \texttt{\tiny{[GM$\langle$P2]}}
                \texttt{СКАЖИ: D успішно на (4,1). Поточні координати: D(4,1), P(3,1), T(3,4), M(6,4), A(1,6), S(6,7), O(7,7). Зараз переміщую T з (3,4) на (1,1), яка порожня.} \\
            }
        }
    }
     \\ \\

    \theutterance \stepcounter{utterance}  
    & & & \multicolumn{2}{p{0.3\linewidth}}{
        \cellcolor[rgb]{0.9,0.9,0.9}{
            \makecell[{{p{\linewidth}}}]{
                \texttt{\tiny{[GM$|$GM]}}
                \texttt{СКАЖИ: D успішно на (4,1). Поточні координати: D(4,1), P(3,1), T(3,4), M(6,4), A(1,6), S(6,7), O(7,7). Зараз переміщую T з (3,4) на (1,1), яка порожня.} \\
            }
        }
    }
    & & \\ \\

    \theutterance \stepcounter{utterance}  
    & & \multicolumn{4}{p{0.6\linewidth}}{
        \cellcolor[rgb]{0.9,0.9,0.9}{
            \makecell[{{p{\linewidth}}}]{
                \texttt{\tiny{[P1$\langle$GM]}}
                \texttt{Ваше повідомлення передано іншому гравцю.} \\
\\ 
\\ 
\texttt{Зараз ви граєте раунд 7 з максимум 28.} \\
\texttt{Разом ви накопичили 6 з 16 штрафів.} \\
\texttt{Інший гравець надіслав це повідомлення:} \\
\\ 
\\ 
\texttt{"D успішно на (4,1). Поточні координати: D(4,1), P(3,1), T(3,4), M(6,4), A(1,6), S(6,7), O(7,7). Зараз переміщую T з (3,4) на (1,1), яка порожня."} \\
\\ 
\\ 
\texttt{Яка ваша наступна команда?} \\
            }
        }
    }
    & \\ \\

    \theutterance \stepcounter{utterance}  
    & \multicolumn{4}{p{0.6\linewidth}}{
        \cellcolor[rgb]{0.9,0.9,0.9}{
            \makecell[{{p{\linewidth}}}]{
                \texttt{\tiny{[P1$\rangle$GM]}}
                \texttt{СКАЖИ: Чудово, рухай T на (1,1). Потім за планом: M на (4,2), A на (1,2), S на (2,7), O на (7,5). Перевіряй порожність перед кожним ходом.} \\
            }
        }
    }
    & & \\ \\

    \theutterance \stepcounter{utterance}  
    & & & \multicolumn{2}{p{0.3\linewidth}}{
        \cellcolor[rgb]{0.9,0.9,0.9}{
            \makecell[{{p{\linewidth}}}]{
                \texttt{\tiny{[GM$|$GM]}}
                \texttt{СКАЖИ: Чудово, рухай T на (1,1). Потім за планом: M на (4,2), A на (1,2), S на (2,7), O на (7,5). Перевіряй порожність перед кожним ходом.} \\
            }
        }
    }
    & & \\ \\

    \theutterance \stepcounter{utterance}  
    & & \multicolumn{4}{p{0.6\linewidth}}{
        \cellcolor[rgb]{0.9,0.9,0.9}{
            \makecell[{{p{\linewidth}}}]{
                \texttt{\tiny{[GM$\rangle$P2]}}
                \texttt{Ваше повідомлення передано іншому гравцю.} \\
\\ 
\\ 
\texttt{Зараз ви граєте раунд 8 з максимум 28.} \\
\texttt{Разом ви накопичили 6 з 16 штрафів.} \\
\texttt{Інший гравець надіслав це повідомлення:} \\
\\ 
\\ 
\texttt{"Чудово, рухай T на (1,1). Потім за планом: M на (4,2), A на (1,2), S на (2,7), O на (7,5). Перевіряй порожність перед кожним ходом."} \\
\\ 
\\ 
\texttt{Яка ваша наступна команда?} \\
            }
        }
    }
    & \\ \\

    \theutterance \stepcounter{utterance}  
    & & & \multicolumn{4}{p{0.6\linewidth}}{
        \cellcolor[rgb]{0.9,0.9,0.9}{
            \makecell[{{p{\linewidth}}}]{
                \texttt{\tiny{[GM$\langle$P2]}}
                \texttt{ПЕРЕСУНЬ: T, (1, 1)} \\
            }
        }
    }
     \\ \\

    \theutterance \stepcounter{utterance}  
    & & & \multicolumn{2}{p{0.3\linewidth}}{
        \cellcolor[rgb]{0.9,0.9,0.9}{
            \makecell[{{p{\linewidth}}}]{
                \texttt{\tiny{[GM$|$GM]}}
                \texttt{ПЕРЕСУНЬ: T, (1, 1)} \\
            }
        }
    }
    & & \\ \\

    \theutterance \stepcounter{utterance}  
    & & & \multicolumn{2}{p{0.3\linewidth}}{
        \cellcolor[rgb]{0.9,0.9,0.9}{
            \makecell[{{p{\linewidth}}}]{
                \texttt{\tiny{[GM$|$GM]}}
                \texttt{T moved to (1, 1)} \\
            }
        }
    }
    & & \\ \\

    \theutterance \stepcounter{utterance}  
    & & \multicolumn{4}{p{0.6\linewidth}}{
        \cellcolor[rgb]{0.9,0.9,0.9}{
            \makecell[{{p{\linewidth}}}]{
                \texttt{\tiny{[P1$\langle$GM]}}
                \texttt{Ваше повідомлення передано іншому гравцю.} \\
\\ 
\\ 
\texttt{Зараз ви граєте раунд 8 з максимум 28.} \\
\texttt{Разом ви накопичили 6 з 16 штрафів.} \\
\texttt{Інший гравець перемістив об'єкт на своїй сітці.} \\
\\ 
\\ 
\texttt{Яка ваша наступна команда?} \\
            }
        }
    }
    & \\ \\

    \theutterance \stepcounter{utterance}  
    & \multicolumn{4}{p{0.6\linewidth}}{
        \cellcolor[rgb]{0.9,0.9,0.9}{
            \makecell[{{p{\linewidth}}}]{
                \texttt{\tiny{[P1$\rangle$GM]}}
                \texttt{СКАЖИ: Як успіхи з T? Які координати зараз?} \\
            }
        }
    }
    & & \\ \\

    \theutterance \stepcounter{utterance}  
    & & & \multicolumn{2}{p{0.3\linewidth}}{
        \cellcolor[rgb]{0.9,0.9,0.9}{
            \makecell[{{p{\linewidth}}}]{
                \texttt{\tiny{[GM$|$GM]}}
                \texttt{СКАЖИ: Як успіхи з T? Які координати зараз?} \\
            }
        }
    }
    & & \\ \\

    \theutterance \stepcounter{utterance}  
    & & \multicolumn{4}{p{0.6\linewidth}}{
        \cellcolor[rgb]{0.9,0.9,0.9}{
            \makecell[{{p{\linewidth}}}]{
                \texttt{\tiny{[GM$\rangle$P2]}}
                \texttt{Об'єкт 'T' успішно переміщено в (1,1). Ваша оновлена сітка виглядає так:} \\
\texttt{```} \\
\texttt{ 1234567} \\
\texttt{╔════╤══╗} \\
\texttt{║T◌PD│◌◌║ 1} \\
\texttt{║◌◌◌◌│◌◌║ 2} \\
\texttt{║◌◌◌◌│◌◌║ 3} \\
\texttt{║◌◌◌◌│M◌║ 4} \\
\texttt{╟──┐◌│◌◌║ 5} \\
\texttt{║A◌├─┼──╢ 6} \\
\texttt{║◌◌│◌│SO║ 7} \\
\texttt{╚══╧═╧══╝} \\
\texttt{```} \\
\texttt{Зараз ви граєте раунд 9 з максимум 28.} \\
\texttt{Разом ви накопичили 6 з 16 штрафів.} \\
\texttt{Інший гравець надіслав це повідомлення:} \\
\\ 
\\ 
\texttt{"Як успіхи з T? Які координати зараз?"} \\
\\ 
\\ 
\texttt{Яка ваша наступна команда?} \\
            }
        }
    }
    & \\ \\

    \theutterance \stepcounter{utterance}  
    & & & \multicolumn{4}{p{0.6\linewidth}}{
        \cellcolor[rgb]{0.9,0.9,0.9}{
            \makecell[{{p{\linewidth}}}]{
                \texttt{\tiny{[GM$\langle$P2]}}
                \texttt{СКАЖИ: T на (1,1). Поточні: T(1,1), P(3,1), D(4,1), M(6,4), A(1,6), S(6,7), O(7,7). Переміщую M з (6,4) на (4,2), перевіряю — (4,2) порожня.} \\
            }
        }
    }
     \\ \\

    \theutterance \stepcounter{utterance}  
    & & & \multicolumn{2}{p{0.3\linewidth}}{
        \cellcolor[rgb]{0.9,0.9,0.9}{
            \makecell[{{p{\linewidth}}}]{
                \texttt{\tiny{[GM$|$GM]}}
                \texttt{СКАЖИ: T на (1,1). Поточні: T(1,1), P(3,1), D(4,1), M(6,4), A(1,6), S(6,7), O(7,7). Переміщую M з (6,4) на (4,2), перевіряю — (4,2) порожня.} \\
            }
        }
    }
    & & \\ \\

    \theutterance \stepcounter{utterance}  
    & & \multicolumn{4}{p{0.6\linewidth}}{
        \cellcolor[rgb]{0.9,0.9,0.9}{
            \makecell[{{p{\linewidth}}}]{
                \texttt{\tiny{[P1$\langle$GM]}}
                \texttt{Ваше повідомлення передано іншому гравцю.} \\
\\ 
\\ 
\texttt{Зараз ви граєте раунд 9 з максимум 28.} \\
\texttt{Разом ви накопичили 6 з 16 штрафів.} \\
\texttt{Інший гравець надіслав це повідомлення:} \\
\\ 
\\ 
\texttt{"T на (1,1). Поточні: T(1,1), P(3,1), D(4,1), M(6,4), A(1,6), S(6,7), O(7,7). Переміщую M з (6,4) на (4,2), перевіряю — (4,2) порожня."} \\
\\ 
\\ 
\texttt{Яка ваша наступна команда?} \\
            }
        }
    }
    & \\ \\

    \theutterance \stepcounter{utterance}  
    & \multicolumn{4}{p{0.6\linewidth}}{
        \cellcolor[rgb]{0.9,0.9,0.9}{
            \makecell[{{p{\linewidth}}}]{
                \texttt{\tiny{[P1$\rangle$GM]}}
                \texttt{СКАЖИ: Добре, рухай M на (4,2). Далі A на (1,2), потім S на (2,7) і O на (7,5). Ти на правильному шляху!} \\
            }
        }
    }
    & & \\ \\

    \theutterance \stepcounter{utterance}  
    & & & \multicolumn{2}{p{0.3\linewidth}}{
        \cellcolor[rgb]{0.9,0.9,0.9}{
            \makecell[{{p{\linewidth}}}]{
                \texttt{\tiny{[GM$|$GM]}}
                \texttt{СКАЖИ: Добре, рухай M на (4,2). Далі A на (1,2), потім S на (2,7) і O на (7,5). Ти на правильному шляху!} \\
            }
        }
    }
    & & \\ \\

    \theutterance \stepcounter{utterance}  
    & & \multicolumn{4}{p{0.6\linewidth}}{
        \cellcolor[rgb]{0.9,0.9,0.9}{
            \makecell[{{p{\linewidth}}}]{
                \texttt{\tiny{[GM$\rangle$P2]}}
                \texttt{Ваше повідомлення передано іншому гравцю.} \\
\\ 
\\ 
\texttt{Зараз ви граєте раунд 10 з максимум 28.} \\
\texttt{Разом ви накопичили 6 з 16 штрафів.} \\
\texttt{Інший гравець надіслав це повідомлення:} \\
\\ 
\\ 
\texttt{"Добре, рухай M на (4,2). Далі A на (1,2), потім S на (2,7) і O на (7,5). Ти на правильному шляху!"} \\
\\ 
\\ 
\texttt{Яка ваша наступна команда?} \\
            }
        }
    }
    & \\ \\

    \theutterance \stepcounter{utterance}  
    & & & \multicolumn{4}{p{0.6\linewidth}}{
        \cellcolor[rgb]{0.9,0.9,0.9}{
            \makecell[{{p{\linewidth}}}]{
                \texttt{\tiny{[GM$\langle$P2]}}
                \texttt{ПЕРЕСУНЬ: M, (4, 2)} \\
            }
        }
    }
     \\ \\

    \theutterance \stepcounter{utterance}  
    & & & \multicolumn{2}{p{0.3\linewidth}}{
        \cellcolor[rgb]{0.9,0.9,0.9}{
            \makecell[{{p{\linewidth}}}]{
                \texttt{\tiny{[GM$|$GM]}}
                \texttt{ПЕРЕСУНЬ: M, (4, 2)} \\
            }
        }
    }
    & & \\ \\

    \theutterance \stepcounter{utterance}  
    & & & \multicolumn{2}{p{0.3\linewidth}}{
        \cellcolor[rgb]{0.9,0.9,0.9}{
            \makecell[{{p{\linewidth}}}]{
                \texttt{\tiny{[GM$|$GM]}}
                \texttt{M moved to (4, 2)} \\
            }
        }
    }
    & & \\ \\

    \theutterance \stepcounter{utterance}  
    & & \multicolumn{4}{p{0.6\linewidth}}{
        \cellcolor[rgb]{0.9,0.9,0.9}{
            \makecell[{{p{\linewidth}}}]{
                \texttt{\tiny{[P1$\langle$GM]}}
                \texttt{Ваше повідомлення передано іншому гравцю.} \\
\\ 
\\ 
\texttt{Зараз ви граєте раунд 10 з максимум 28.} \\
\texttt{Разом ви накопичили 6 з 16 штрафів.} \\
\texttt{Інший гравець перемістив об'єкт на своїй сітці.} \\
\\ 
\\ 
\texttt{Яка ваша наступна команда?} \\
            }
        }
    }
    & \\ \\

    \theutterance \stepcounter{utterance}  
    & \multicolumn{4}{p{0.6\linewidth}}{
        \cellcolor[rgb]{0.9,0.9,0.9}{
            \makecell[{{p{\linewidth}}}]{
                \texttt{\tiny{[P1$\rangle$GM]}}
                \texttt{СКАЖИ: Як пройшло з M? Які поточні координати?} \\
            }
        }
    }
    & & \\ \\

    \theutterance \stepcounter{utterance}  
    & & & \multicolumn{2}{p{0.3\linewidth}}{
        \cellcolor[rgb]{0.9,0.9,0.9}{
            \makecell[{{p{\linewidth}}}]{
                \texttt{\tiny{[GM$|$GM]}}
                \texttt{СКАЖИ: Як пройшло з M? Які поточні координати?} \\
            }
        }
    }
    & & \\ \\

    \theutterance \stepcounter{utterance}  
    & & \multicolumn{4}{p{0.6\linewidth}}{
        \cellcolor[rgb]{0.9,0.9,0.9}{
            \makecell[{{p{\linewidth}}}]{
                \texttt{\tiny{[GM$\rangle$P2]}}
                \texttt{Об'єкт 'M' успішно переміщено в (4,2). Ваша оновлена сітка виглядає так:} \\
\texttt{```} \\
\texttt{ 1234567} \\
\texttt{╔════╤══╗} \\
\texttt{║T◌PD│◌◌║ 1} \\
\texttt{║◌◌◌M│◌◌║ 2} \\
\texttt{║◌◌◌◌│◌◌║ 3} \\
\texttt{║◌◌◌◌│◌◌║ 4} \\
\texttt{╟──┐◌│◌◌║ 5} \\
\texttt{║A◌├─┼──╢ 6} \\
\texttt{║◌◌│◌│SO║ 7} \\
\texttt{╚══╧═╧══╝} \\
\texttt{```} \\
\texttt{Зараз ви граєте раунд 11 з максимум 28.} \\
\texttt{Разом ви накопичили 6 з 16 штрафів.} \\
\texttt{Інший гравець надіслав це повідомлення:} \\
\\ 
\\ 
\texttt{"Як пройшло з M? Які поточні координати?"} \\
\\ 
\\ 
\texttt{Яка ваша наступна команда?} \\
            }
        }
    }
    & \\ \\

    \theutterance \stepcounter{utterance}  
    & & & \multicolumn{4}{p{0.6\linewidth}}{
        \cellcolor[rgb]{0.9,0.9,0.9}{
            \makecell[{{p{\linewidth}}}]{
                \texttt{\tiny{[GM$\langle$P2]}}
                \texttt{СКАЖИ: M на (4,2). Поточні: T(1,1), P(3,1), D(4,1), M(4,2), A(1,6), S(6,7), O(7,7). Переміщую A з (1,6) на (1,2), вона порожня.} \\
            }
        }
    }
     \\ \\

    \theutterance \stepcounter{utterance}  
    & & & \multicolumn{2}{p{0.3\linewidth}}{
        \cellcolor[rgb]{0.9,0.9,0.9}{
            \makecell[{{p{\linewidth}}}]{
                \texttt{\tiny{[GM$|$GM]}}
                \texttt{СКАЖИ: M на (4,2). Поточні: T(1,1), P(3,1), D(4,1), M(4,2), A(1,6), S(6,7), O(7,7). Переміщую A з (1,6) на (1,2), вона порожня.} \\
            }
        }
    }
    & & \\ \\

    \theutterance \stepcounter{utterance}  
    & & \multicolumn{4}{p{0.6\linewidth}}{
        \cellcolor[rgb]{0.9,0.9,0.9}{
            \makecell[{{p{\linewidth}}}]{
                \texttt{\tiny{[P1$\langle$GM]}}
                \texttt{Ваше повідомлення передано іншому гравцю.} \\
\\ 
\\ 
\texttt{Зараз ви граєте раунд 11 з максимум 28.} \\
\texttt{Разом ви накопичили 6 з 16 штрафів.} \\
\texttt{Інший гравець надіслав це повідомлення:} \\
\\ 
\\ 
\texttt{"M на (4,2). Поточні: T(1,1), P(3,1), D(4,1), M(4,2), A(1,6), S(6,7), O(7,7). Переміщую A з (1,6) на (1,2), вона порожня."} \\
\\ 
\\ 
\texttt{Яка ваша наступна команда?} \\
            }
        }
    }
    & \\ \\

    \theutterance \stepcounter{utterance}  
    & \multicolumn{4}{p{0.6\linewidth}}{
        \cellcolor[rgb]{0.9,0.9,0.9}{
            \makecell[{{p{\linewidth}}}]{
                \texttt{\tiny{[P1$\rangle$GM]}}
                \texttt{СКАЖИ: Рухай A на (1,2). Залишилися тільки S на (2,7) та O на (7,5). Майже готово!} \\
            }
        }
    }
    & & \\ \\

    \theutterance \stepcounter{utterance}  
    & & & \multicolumn{2}{p{0.3\linewidth}}{
        \cellcolor[rgb]{0.9,0.9,0.9}{
            \makecell[{{p{\linewidth}}}]{
                \texttt{\tiny{[GM$|$GM]}}
                \texttt{СКАЖИ: Рухай A на (1,2). Залишилися тільки S на (2,7) та O на (7,5). Майже готово!} \\
            }
        }
    }
    & & \\ \\

    \theutterance \stepcounter{utterance}  
    & & \multicolumn{4}{p{0.6\linewidth}}{
        \cellcolor[rgb]{0.9,0.9,0.9}{
            \makecell[{{p{\linewidth}}}]{
                \texttt{\tiny{[GM$\rangle$P2]}}
                \texttt{Ваше повідомлення передано іншому гравцю.} \\
\\ 
\\ 
\texttt{Зараз ви граєте раунд 12 з максимум 28.} \\
\texttt{Разом ви накопичили 6 з 16 штрафів.} \\
\texttt{Інший гравець надіслав це повідомлення:} \\
\\ 
\\ 
\texttt{"Рухай A на (1,2). Залишилися тільки S на (2,7) та O на (7,5). Майже готово!"} \\
\\ 
\\ 
\texttt{Яка ваша наступна команда?} \\
            }
        }
    }
    & \\ \\

    \theutterance \stepcounter{utterance}  
    & & & \multicolumn{4}{p{0.6\linewidth}}{
        \cellcolor[rgb]{0.9,0.9,0.9}{
            \makecell[{{p{\linewidth}}}]{
                \texttt{\tiny{[GM$\langle$P2]}}
                \texttt{ПЕРЕСУНЬ: A, (1, 2)} \\
            }
        }
    }
     \\ \\

    \theutterance \stepcounter{utterance}  
    & & & \multicolumn{2}{p{0.3\linewidth}}{
        \cellcolor[rgb]{0.9,0.9,0.9}{
            \makecell[{{p{\linewidth}}}]{
                \texttt{\tiny{[GM$|$GM]}}
                \texttt{ПЕРЕСУНЬ: A, (1, 2)} \\
            }
        }
    }
    & & \\ \\

    \theutterance \stepcounter{utterance}  
    & & & \multicolumn{2}{p{0.3\linewidth}}{
        \cellcolor[rgb]{0.9,0.9,0.9}{
            \makecell[{{p{\linewidth}}}]{
                \texttt{\tiny{[GM$|$GM]}}
                \texttt{A moved to (1, 2)} \\
            }
        }
    }
    & & \\ \\

    \theutterance \stepcounter{utterance}  
    & & \multicolumn{4}{p{0.6\linewidth}}{
        \cellcolor[rgb]{0.9,0.9,0.9}{
            \makecell[{{p{\linewidth}}}]{
                \texttt{\tiny{[P1$\langle$GM]}}
                \texttt{Ваше повідомлення передано іншому гравцю.} \\
\\ 
\\ 
\texttt{Зараз ви граєте раунд 12 з максимум 28.} \\
\texttt{Разом ви накопичили 6 з 16 штрафів.} \\
\texttt{Інший гравець перемістив об'єкт на своїй сітці.} \\
\\ 
\\ 
\texttt{Яка ваша наступна команда?} \\
            }
        }
    }
    & \\ \\

    \theutterance \stepcounter{utterance}  
    & \multicolumn{4}{p{0.6\linewidth}}{
        \cellcolor[rgb]{0.9,0.9,0.9}{
            \makecell[{{p{\linewidth}}}]{
                \texttt{\tiny{[P1$\rangle$GM]}}
                \texttt{СКАЖИ: A на місці? Які поточні координати?} \\
            }
        }
    }
    & & \\ \\

    \theutterance \stepcounter{utterance}  
    & & & \multicolumn{2}{p{0.3\linewidth}}{
        \cellcolor[rgb]{0.9,0.9,0.9}{
            \makecell[{{p{\linewidth}}}]{
                \texttt{\tiny{[GM$|$GM]}}
                \texttt{СКАЖИ: A на місці? Які поточні координати?} \\
            }
        }
    }
    & & \\ \\

    \theutterance \stepcounter{utterance}  
    & & \multicolumn{4}{p{0.6\linewidth}}{
        \cellcolor[rgb]{0.9,0.9,0.9}{
            \makecell[{{p{\linewidth}}}]{
                \texttt{\tiny{[GM$\rangle$P2]}}
                \texttt{Об'єкт 'A' успішно переміщено в (1,2). Ваша оновлена сітка виглядає так:} \\
\texttt{```} \\
\texttt{ 1234567} \\
\texttt{╔════╤══╗} \\
\texttt{║T◌PD│◌◌║ 1} \\
\texttt{║A◌◌M│◌◌║ 2} \\
\texttt{║◌◌◌◌│◌◌║ 3} \\
\texttt{║◌◌◌◌│◌◌║ 4} \\
\texttt{╟──┐◌│◌◌║ 5} \\
\texttt{║◌◌├─┼──╢ 6} \\
\texttt{║◌◌│◌│SO║ 7} \\
\texttt{╚══╧═╧══╝} \\
\texttt{```} \\
\texttt{Зараз ви граєте раунд 13 з максимум 28.} \\
\texttt{Разом ви накопичили 6 з 16 штрафів.} \\
\texttt{Інший гравець надіслав це повідомлення:} \\
\\ 
\\ 
\texttt{"A на місці? Які поточні координати?"} \\
\\ 
\\ 
\texttt{Яка ваша наступна команда?} \\
            }
        }
    }
    & \\ \\

    \theutterance \stepcounter{utterance}  
    & & & \multicolumn{4}{p{0.6\linewidth}}{
        \cellcolor[rgb]{0.9,0.9,0.9}{
            \makecell[{{p{\linewidth}}}]{
                \texttt{\tiny{[GM$\langle$P2]}}
                \texttt{СКАЖИ: A на (1,2). Поточні: T(1,1), P(3,1), D(4,1), A(1,2), M(4,2), S(6,7), O(7,7). Переміщую S з (6,7) на (2,7), перевіряю — (2,7) порожня.} \\
            }
        }
    }
     \\ \\

    \theutterance \stepcounter{utterance}  
    & & & \multicolumn{2}{p{0.3\linewidth}}{
        \cellcolor[rgb]{0.9,0.9,0.9}{
            \makecell[{{p{\linewidth}}}]{
                \texttt{\tiny{[GM$|$GM]}}
                \texttt{СКАЖИ: A на (1,2). Поточні: T(1,1), P(3,1), D(4,1), A(1,2), M(4,2), S(6,7), O(7,7). Переміщую S з (6,7) на (2,7), перевіряю — (2,7) порожня.} \\
            }
        }
    }
    & & \\ \\

    \theutterance \stepcounter{utterance}  
    & & \multicolumn{4}{p{0.6\linewidth}}{
        \cellcolor[rgb]{0.9,0.9,0.9}{
            \makecell[{{p{\linewidth}}}]{
                \texttt{\tiny{[P1$\langle$GM]}}
                \texttt{Ваше повідомлення передано іншому гравцю.} \\
\\ 
\\ 
\texttt{Зараз ви граєте раунд 13 з максимум 28.} \\
\texttt{Разом ви накопичили 6 з 16 штрафів.} \\
\texttt{Інший гравець надіслав це повідомлення:} \\
\\ 
\\ 
\texttt{"A на (1,2). Поточні: T(1,1), P(3,1), D(4,1), A(1,2), M(4,2), S(6,7), O(7,7). Переміщую S з (6,7) на (2,7), перевіряю — (2,7) порожня."} \\
\\ 
\\ 
\texttt{Яка ваша наступна команда?} \\
            }
        }
    }
    & \\ \\

    \theutterance \stepcounter{utterance}  
    & \multicolumn{4}{p{0.6\linewidth}}{
        \cellcolor[rgb]{0.9,0.9,0.9}{
            \makecell[{{p{\linewidth}}}]{
                \texttt{\tiny{[P1$\rangle$GM]}}
                \texttt{СКАЖИ: Рухай S на (2,7). Останнім буде O на (7,5). Ми майже виграли!} \\
            }
        }
    }
    & & \\ \\

    \theutterance \stepcounter{utterance}  
    & & & \multicolumn{2}{p{0.3\linewidth}}{
        \cellcolor[rgb]{0.9,0.9,0.9}{
            \makecell[{{p{\linewidth}}}]{
                \texttt{\tiny{[GM$|$GM]}}
                \texttt{СКАЖИ: Рухай S на (2,7). Останнім буде O на (7,5). Ми майже виграли!} \\
            }
        }
    }
    & & \\ \\

    \theutterance \stepcounter{utterance}  
    & & \multicolumn{4}{p{0.6\linewidth}}{
        \cellcolor[rgb]{0.9,0.9,0.9}{
            \makecell[{{p{\linewidth}}}]{
                \texttt{\tiny{[GM$\rangle$P2]}}
                \texttt{Ваше повідомлення передано іншому гравцю.} \\
\\ 
\\ 
\texttt{Зараз ви граєте раунд 14 з максимум 28.} \\
\texttt{Разом ви накопичили 6 з 16 штрафів.} \\
\texttt{Інший гравець надіслав це повідомлення:} \\
\\ 
\\ 
\texttt{"Рухай S на (2,7). Останнім буде O на (7,5). Ми майже виграли!"} \\
\\ 
\\ 
\texttt{Яка ваша наступна команда?} \\
            }
        }
    }
    & \\ \\

    \theutterance \stepcounter{utterance}  
    & & & \multicolumn{4}{p{0.6\linewidth}}{
        \cellcolor[rgb]{0.9,0.9,0.9}{
            \makecell[{{p{\linewidth}}}]{
                \texttt{\tiny{[GM$\langle$P2]}}
                \texttt{ПЕРЕСУНЬ: S, (2, 7)} \\
            }
        }
    }
     \\ \\

    \theutterance \stepcounter{utterance}  
    & & & \multicolumn{2}{p{0.3\linewidth}}{
        \cellcolor[rgb]{0.9,0.9,0.9}{
            \makecell[{{p{\linewidth}}}]{
                \texttt{\tiny{[GM$|$GM]}}
                \texttt{ПЕРЕСУНЬ: S, (2, 7)} \\
            }
        }
    }
    & & \\ \\

    \theutterance \stepcounter{utterance}  
    & & & \multicolumn{2}{p{0.3\linewidth}}{
        \cellcolor[rgb]{0.9,0.9,0.9}{
            \makecell[{{p{\linewidth}}}]{
                \texttt{\tiny{[GM$|$GM]}}
                \texttt{S moved to (2, 7)} \\
            }
        }
    }
    & & \\ \\

    \theutterance \stepcounter{utterance}  
    & & \multicolumn{4}{p{0.6\linewidth}}{
        \cellcolor[rgb]{0.9,0.9,0.9}{
            \makecell[{{p{\linewidth}}}]{
                \texttt{\tiny{[P1$\langle$GM]}}
                \texttt{Ваше повідомлення передано іншому гравцю.} \\
\\ 
\\ 
\texttt{Зараз ви граєте раунд 14 з максимум 28.} \\
\texttt{Разом ви накопичили 6 з 16 штрафів.} \\
\texttt{Інший гравець перемістив об'єкт на своїй сітці.} \\
\\ 
\\ 
\texttt{Яка ваша наступна команда?} \\
            }
        }
    }
    & \\ \\

    \theutterance \stepcounter{utterance}  
    & \multicolumn{4}{p{0.6\linewidth}}{
        \cellcolor[rgb]{0.9,0.9,0.9}{
            \makecell[{{p{\linewidth}}}]{
                \texttt{\tiny{[P1$\rangle$GM]}}
                \texttt{СКАЖИ: S на місці? Які координати зараз?} \\
            }
        }
    }
    & & \\ \\

    \theutterance \stepcounter{utterance}  
    & & & \multicolumn{2}{p{0.3\linewidth}}{
        \cellcolor[rgb]{0.9,0.9,0.9}{
            \makecell[{{p{\linewidth}}}]{
                \texttt{\tiny{[GM$|$GM]}}
                \texttt{СКАЖИ: S на місці? Які координати зараз?} \\
            }
        }
    }
    & & \\ \\

    \theutterance \stepcounter{utterance}  
    & & \multicolumn{4}{p{0.6\linewidth}}{
        \cellcolor[rgb]{0.9,0.9,0.9}{
            \makecell[{{p{\linewidth}}}]{
                \texttt{\tiny{[GM$\rangle$P2]}}
                \texttt{Об'єкт 'S' успішно переміщено в (2,7). Ваша оновлена сітка виглядає так:} \\
\texttt{```} \\
\texttt{ 1234567} \\
\texttt{╔════╤══╗} \\
\texttt{║T◌PD│◌◌║ 1} \\
\texttt{║A◌◌M│◌◌║ 2} \\
\texttt{║◌◌◌◌│◌◌║ 3} \\
\texttt{║◌◌◌◌│◌◌║ 4} \\
\texttt{╟──┐◌│◌◌║ 5} \\
\texttt{║◌◌├─┼──╢ 6} \\
\texttt{║◌S│◌│◌O║ 7} \\
\texttt{╚══╧═╧══╝} \\
\texttt{```} \\
\texttt{Зараз ви граєте раунд 15 з максимум 28.} \\
\texttt{Разом ви накопичили 6 з 16 штрафів.} \\
\texttt{Інший гравець надіслав це повідомлення:} \\
\\ 
\\ 
\texttt{"S на місці? Які координати зараз?"} \\
\\ 
\\ 
\texttt{Яка ваша наступна команда?} \\
            }
        }
    }
    & \\ \\

    \theutterance \stepcounter{utterance}  
    & & & \multicolumn{4}{p{0.6\linewidth}}{
        \cellcolor[rgb]{0.9,0.9,0.9}{
            \makecell[{{p{\linewidth}}}]{
                \texttt{\tiny{[GM$\langle$P2]}}
                \texttt{СКАЖИ: S на (2,7). Поточні: T(1,1), P(3,1), D(4,1), A(1,2), M(4,2), S(2,7), O(7,7). Переміщую O з (7,7) на (7,5), перевіряю — (7,5) порожня.} \\
            }
        }
    }
     \\ \\

    \theutterance \stepcounter{utterance}  
    & & & \multicolumn{2}{p{0.3\linewidth}}{
        \cellcolor[rgb]{0.9,0.9,0.9}{
            \makecell[{{p{\linewidth}}}]{
                \texttt{\tiny{[GM$|$GM]}}
                \texttt{СКАЖИ: S на (2,7). Поточні: T(1,1), P(3,1), D(4,1), A(1,2), M(4,2), S(2,7), O(7,7). Переміщую O з (7,7) на (7,5), перевіряю — (7,5) порожня.} \\
            }
        }
    }
    & & \\ \\

    \theutterance \stepcounter{utterance}  
    & & \multicolumn{4}{p{0.6\linewidth}}{
        \cellcolor[rgb]{0.9,0.9,0.9}{
            \makecell[{{p{\linewidth}}}]{
                \texttt{\tiny{[P1$\langle$GM]}}
                \texttt{Ваше повідомлення передано іншому гравцю.} \\
\\ 
\\ 
\texttt{Зараз ви граєте раунд 15 з максимум 28.} \\
\texttt{Разом ви накопичили 6 з 16 штрафів.} \\
\texttt{Інший гравець надіслав це повідомлення:} \\
\\ 
\\ 
\texttt{"S на (2,7). Поточні: T(1,1), P(3,1), D(4,1), A(1,2), M(4,2), S(2,7), O(7,7). Переміщую O з (7,7) на (7,5), перевіряю — (7,5) порожня."} \\
\\ 
\\ 
\texttt{Яка ваша наступна команда?} \\
            }
        }
    }
    & \\ \\

    \theutterance \stepcounter{utterance}  
    & \multicolumn{4}{p{0.6\linewidth}}{
        \cellcolor[rgb]{0.9,0.9,0.9}{
            \makecell[{{p{\linewidth}}}]{
                \texttt{\tiny{[P1$\rangle$GM]}}
                \texttt{СКАЖИ: Рухай O на (7,5). Це останній об'єкт! Після цього всі координати збігатимуться з моїми: T(1,1), P(3,1), D(4,1), A(1,2), M(4,2), S(2,7), O(7,5). Завершуємо гру!} \\
            }
        }
    }
    & & \\ \\

    \theutterance \stepcounter{utterance}  
    & & & \multicolumn{2}{p{0.3\linewidth}}{
        \cellcolor[rgb]{0.9,0.9,0.9}{
            \makecell[{{p{\linewidth}}}]{
                \texttt{\tiny{[GM$|$GM]}}
                \texttt{СКАЖИ: Рухай O на (7,5). Це останній об'єкт! Після цього всі координати збігатимуться з моїми: T(1,1), P(3,1), D(4,1), A(1,2), M(4,2), S(2,7), O(7,5). Завершуємо гру!} \\
            }
        }
    }
    & & \\ \\

    \theutterance \stepcounter{utterance}  
    & & \multicolumn{4}{p{0.6\linewidth}}{
        \cellcolor[rgb]{0.9,0.9,0.9}{
            \makecell[{{p{\linewidth}}}]{
                \texttt{\tiny{[GM$\rangle$P2]}}
                \texttt{Ваше повідомлення передано іншому гравцю.} \\
\\ 
\\ 
\texttt{Зараз ви граєте раунд 16 з максимум 28.} \\
\texttt{Разом ви накопичили 6 з 16 штрафів.} \\
\texttt{Інший гравець надіслав це повідомлення:} \\
\\ 
\\ 
\texttt{"Рухай O на (7,5). Це останній об'єкт! Після цього всі координати збігатимуться з моїми: T(1,1), P(3,1), D(4,1), A(1,2), M(4,2), S(2,7), O(7,5). Завершуємо гру!"} \\
\\ 
\\ 
\texttt{Яка ваша наступна команда?} \\
            }
        }
    }
    & \\ \\

    \theutterance \stepcounter{utterance}  
    & & & \multicolumn{4}{p{0.6\linewidth}}{
        \cellcolor[rgb]{0.9,0.9,0.9}{
            \makecell[{{p{\linewidth}}}]{
                \texttt{\tiny{[GM$\langle$P2]}}
                \texttt{ПЕРЕСУНЬ: O, (7, 5)} \\
            }
        }
    }
     \\ \\

    \theutterance \stepcounter{utterance}  
    & & & \multicolumn{2}{p{0.3\linewidth}}{
        \cellcolor[rgb]{0.9,0.9,0.9}{
            \makecell[{{p{\linewidth}}}]{
                \texttt{\tiny{[GM$|$GM]}}
                \texttt{ПЕРЕСУНЬ: O, (7, 5)} \\
            }
        }
    }
    & & \\ \\

    \theutterance \stepcounter{utterance}  
    & & & \multicolumn{2}{p{0.3\linewidth}}{
        \cellcolor[rgb]{0.9,0.9,0.9}{
            \makecell[{{p{\linewidth}}}]{
                \texttt{\tiny{[GM$|$GM]}}
                \texttt{O moved to (7, 5)} \\
            }
        }
    }
    & & \\ \\

    \theutterance \stepcounter{utterance}  
    & & \multicolumn{4}{p{0.6\linewidth}}{
        \cellcolor[rgb]{0.9,0.9,0.9}{
            \makecell[{{p{\linewidth}}}]{
                \texttt{\tiny{[P1$\langle$GM]}}
                \texttt{Ваше повідомлення передано іншому гравцю.} \\
\\ 
\\ 
\texttt{Зараз ви граєте раунд 16 з максимум 28.} \\
\texttt{Разом ви накопичили 6 з 16 штрафів.} \\
\texttt{Інший гравець перемістив об'єкт на своїй сітці.} \\
\\ 
\\ 
\texttt{Яка ваша наступна команда?} \\
            }
        }
    }
    & \\ \\

    \theutterance \stepcounter{utterance}  
    & \multicolumn{4}{p{0.6\linewidth}}{
        \cellcolor[rgb]{0.9,0.9,0.9}{
            \makecell[{{p{\linewidth}}}]{
                \texttt{\tiny{[P1$\rangle$GM]}}
                \texttt{СКАЖИ: O на місці? Всі координати збігаються? Якщо так, завершено?} \\
            }
        }
    }
    & & \\ \\

    \theutterance \stepcounter{utterance}  
    & & & \multicolumn{2}{p{0.3\linewidth}}{
        \cellcolor[rgb]{0.9,0.9,0.9}{
            \makecell[{{p{\linewidth}}}]{
                \texttt{\tiny{[GM$|$GM]}}
                \texttt{СКАЖИ: O на місці? Всі координати збігаються? Якщо так, завершено?} \\
            }
        }
    }
    & & \\ \\

    \theutterance \stepcounter{utterance}  
    & & \multicolumn{4}{p{0.6\linewidth}}{
        \cellcolor[rgb]{0.9,0.9,0.9}{
            \makecell[{{p{\linewidth}}}]{
                \texttt{\tiny{[GM$\rangle$P2]}}
                \texttt{Об'єкт 'O' успішно переміщено в (7,5). Ваша оновлена сітка виглядає так:} \\
\texttt{```} \\
\texttt{ 1234567} \\
\texttt{╔════╤══╗} \\
\texttt{║T◌PD│◌◌║ 1} \\
\texttt{║A◌◌M│◌◌║ 2} \\
\texttt{║◌◌◌◌│◌◌║ 3} \\
\texttt{║◌◌◌◌│◌◌║ 4} \\
\texttt{╟──┐◌│◌O║ 5} \\
\texttt{║◌◌├─┼──╢ 6} \\
\texttt{║◌S│◌│◌◌║ 7} \\
\texttt{╚══╧═╧══╝} \\
\texttt{```} \\
\texttt{Зараз ви граєте раунд 17 з максимум 28.} \\
\texttt{Разом ви накопичили 6 з 16 штрафів.} \\
\texttt{Інший гравець надіслав це повідомлення:} \\
\\ 
\\ 
\texttt{"O на місці? Всі координати збігаються? Якщо так, завершено?"} \\
\\ 
\\ 
\texttt{Яка ваша наступна команда?} \\
            }
        }
    }
    & \\ \\

    \theutterance \stepcounter{utterance}  
    & & & \multicolumn{4}{p{0.6\linewidth}}{
        \cellcolor[rgb]{0.9,0.9,0.9}{
            \makecell[{{p{\linewidth}}}]{
                \texttt{\tiny{[GM$\langle$P2]}}
                \texttt{СКАЖИ: завершено!} \\
            }
        }
    }
     \\ \\

    \theutterance \stepcounter{utterance}  
    & & & \multicolumn{2}{p{0.3\linewidth}}{
        \cellcolor[rgb]{0.9,0.9,0.9}{
            \makecell[{{p{\linewidth}}}]{
                \texttt{\tiny{[GM$|$GM]}}
                \texttt{СКАЖИ: завершено!} \\
            }
        }
    }
    & & \\ \\

    \theutterance \stepcounter{utterance}  
    & & & \multicolumn{2}{p{0.3\linewidth}}{
        \cellcolor[rgb]{0.9,0.9,0.9}{
            \makecell[{{p{\linewidth}}}]{
                \texttt{\tiny{[GM$|$GM]}}
                \texttt{true} \\
            }
        }
    }
    & & \\ \\

    \theutterance \stepcounter{utterance}  
    & & & \multicolumn{2}{p{0.3\linewidth}}{
        \cellcolor[rgb]{0.9,0.9,0.9}{
            \makecell[{{p{\linewidth}}}]{
                \texttt{\tiny{[GM$|$GM]}}
                \texttt{Initial states:} \\
\texttt{```} \\
\texttt{Player 1:} \\
\texttt{ 1234567} \\
\texttt{╔════╤══╗} \\
\texttt{║T◌PD│◌◌║ 1} \\
\texttt{║A◌◌M│◌◌║ 2} \\
\texttt{║◌◌◌◌│◌◌║ 3} \\
\texttt{║◌◌◌◌│◌◌║ 4} \\
\texttt{╟──┐◌│◌O║ 5} \\
\texttt{║◌◌├─┼──╢ 6} \\
\texttt{║◌S│◌│◌◌║ 7} \\
\texttt{╚══╧═╧══╝} \\
\\ 
\\ 
\texttt{Player 2:} \\
\texttt{ 1234567} \\
\texttt{╔════╤══╗} \\
\texttt{║D◌◌P│◌◌║ 1} \\
\texttt{║◌◌◌◌│◌◌║ 2} \\
\texttt{║◌◌◌◌│◌◌║ 3} \\
\texttt{║◌◌T◌│M◌║ 4} \\
\texttt{╟──┐◌│◌◌║ 5} \\
\texttt{║A◌├─┼──╢ 6} \\
\texttt{║◌◌│◌│SO║ 7} \\
\texttt{╚══╧═╧══╝} \\
\\ 
\texttt{```} \\
            }
        }
    }
    & & \\ \\

    \theutterance \stepcounter{utterance}  
    & & & \multicolumn{2}{p{0.3\linewidth}}{
        \cellcolor[rgb]{0.9,0.9,0.9}{
            \makecell[{{p{\linewidth}}}]{
                \texttt{\tiny{[GM$|$GM]}}
                \texttt{End states:} \\
\texttt{```} \\
\texttt{Player 1:} \\
\texttt{ 1234567} \\
\texttt{╔════╤══╗} \\
\texttt{║T◌PD│◌◌║ 1} \\
\texttt{║A◌◌M│◌◌║ 2} \\
\texttt{║◌◌◌◌│◌◌║ 3} \\
\texttt{║◌◌◌◌│◌◌║ 4} \\
\texttt{╟──┐◌│◌O║ 5} \\
\texttt{║◌◌├─┼──╢ 6} \\
\texttt{║◌S│◌│◌◌║ 7} \\
\texttt{╚══╧═╧══╝} \\
\\ 
\\ 
\texttt{Player 2:} \\
\texttt{ 1234567} \\
\texttt{╔════╤══╗} \\
\texttt{║T◌PD│◌◌║ 1} \\
\texttt{║A◌◌M│◌◌║ 2} \\
\texttt{║◌◌◌◌│◌◌║ 3} \\
\texttt{║◌◌◌◌│◌◌║ 4} \\
\texttt{╟──┐◌│◌O║ 5} \\
\texttt{║◌◌├─┼──╢ 6} \\
\texttt{║◌S│◌│◌◌║ 7} \\
\texttt{╚══╧═╧══╝} \\
\\ 
\texttt{```} \\
            }
        }
    }
    & & \\ \\

    \theutterance \stepcounter{utterance}  
    & & & \multicolumn{2}{p{0.3\linewidth}}{
        \cellcolor[rgb]{0.9,0.9,0.9}{
            \makecell[{{p{\linewidth}}}]{
                \texttt{\tiny{[GM$|$GM]}}
                \texttt{* success: True} \\
\texttt{* lose: False} \\
\texttt{* aborted: False} \\
\texttt{{-}{-}{-}{-}{-}{-}{-}} \\
\texttt{* Shifts: 6.00} \\
\texttt{* Max Shifts: 12.00} \\
\texttt{* Min Shifts: 6.00} \\
\texttt{* End Distance Sum: 0.00} \\
\texttt{* Init Distance Sum: 20.43} \\
\texttt{* Expected Distance Sum: 29.33} \\
\texttt{* Penalties: 6.00} \\
\texttt{* Max Penalties: 16.00} \\
\texttt{* Rounds: 17.00} \\
\texttt{* Max Rounds: 28.00} \\
\texttt{* Object Count: 7.00} \\
            }
        }
    }
    & & \\ \\

\end{supertabular}
}

\end{document}
